%%%%%%%%%%%%%%%%%%%%%%%%%%%%%%%%%%%%%%%%%%%%%%%%%%%%%%%%%%%%
% 11.10 MATHEMATICAL SIMULATION
% The Composition of Advanced Mathematics
%%%%%%%%%%%%%%%%%%%%%%%%%%%%%%%%%%%%%%%%%%%%%%%%%%%%%%%%%%%%

\section{Mathematical Simulation}
\label{sec:mathematical-simulation}

\prerequisite{
Mathematical modeling, probability, differential equations, numerical
analysis, numerical ODEs, numerical PDEs, Monte Carlo methods, optimization,
statistics, and scientific computing.
}

\learningobjectives{
By the end of this section, the reader should be able to:
\begin{itemize}
    \item explain the role of mathematical simulation;
    \item distinguish models from simulations;
    \item distinguish deterministic and stochastic simulation;
    \item distinguish static and dynamic simulation;
    \item identify states, parameters, inputs, and outputs;
    \item construct simulation state-update equations;
    \item understand continuous-time and discrete-time models;
    \item recognize event-driven simulation;
    \item distinguish time-driven and event-driven simulation;
    \item recognize agent-based simulation;
    \item understand compartment models;
    \item recognize network-based simulation;
    \item formulate simulation experiments;
    \item select initial conditions and parameter values;
    \item understand transient and steady-state behavior;
    \item recognize equilibration and warm-up periods;
    \item understand simulation horizons;
    \item understand time-step selection;
    \item distinguish numerical instability from physical instability;
    \item understand conservation diagnostics;
    \item recognize invariant quantities;
    \item construct Monte Carlo simulation studies;
    \item estimate simulation uncertainty;
    \item recognize stochastic trajectories;
    \item understand replication and ensemble averaging;
    \item distinguish aleatoric and parametric uncertainty;
    \item understand sensitivity analysis;
    \item recognize local and global sensitivity methods;
    \item understand parameter sweeps;
    \item recognize scenario analysis;
    \item understand calibration;
    \item distinguish calibration from validation;
    \item understand parameter estimation from simulation outputs;
    \item recognize inverse problems;
    \item understand objective functions for calibration;
    \item understand identifiability conceptually;
    \item recognize structural and practical identifiability;
    \item understand verification and validation;
    \item recognize benchmark problems;
    \item understand convergence studies;
    \item construct manufactured or synthetic-data tests;
    \item recognize uncertainty quantification;
    \item understand output distributions;
    \item estimate probabilities and quantiles from simulation;
    \item recognize rare-event simulation;
    \item understand surrogate models conceptually;
    \item recognize reduced-order simulation;
    \item understand response surfaces;
    \item recognize simulation-based optimization;
    \item understand digital experiments;
    \item recognize reproducible simulation workflows;
    \item understand checkpointing and restart;
    \item recognize parallel ensemble simulation;
    \item understand visualization of dynamic systems;
    \item distinguish visual plausibility from mathematical verification;
    \item connect simulation with biology, physics, engineering, finance,
          operations research, epidemiology, environmental science, and
          applied mathematics.
\end{itemize}
}

%%%%%%%%%%%%%%%%%%%%%%%%%%%%%%%%%%%%%%%%%%%%%%%%%%%%%%%%%%%%
\subsection{What Is Mathematical Simulation?}
\label{subsec:what-is-mathematical-simulation}
%%%%%%%%%%%%%%%%%%%%%%%%%%%%%%%%%%%%%%%%%%%%%%%%%%%%%%%%%%%%

A mathematical simulation uses a mathematical model and a computational
procedure to imitate the behavior of a system.

A typical pipeline is

\[
\boxed{
\text{real system}
\longrightarrow
\text{mathematical model}
\longrightarrow
\text{numerical algorithm}
\longrightarrow
\text{simulation output}.
}
\]

The simulation is not the real system itself.

It is the computed behavior of a mathematical representation.

%%%%%%%%%%%%%%%%%%%%%%%%%%%%%%%%%%%%%%%%%%%%%%%%%%%%%%%%%%%%
\subsection{Model Versus Simulation}
\label{subsec:model-versus-simulation}
%%%%%%%%%%%%%%%%%%%%%%%%%%%%%%%%%%%%%%%%%%%%%%%%%%%%%%%%%%%%

A mathematical model specifies relationships among variables.

For example,

\[
\boxed{
\frac{dP}{dt}
=
rP
\left(
1-\frac{P}{K}
\right)
}
\]

is a logistic population model.

A simulation chooses:

\[
\boxed{
P(0),
}
\]

\[
\boxed{
r,
}
\]

\[
\boxed{
K,
}
\]

a numerical method, and a time interval, then computes an approximate
trajectory.

Thus the same model can produce many different simulations.

%%%%%%%%%%%%%%%%%%%%%%%%%%%%%%%%%%%%%%%%%%%%%%%%%%%%%%%%%%%%
\subsection{Components of a Simulation Model}
\label{subsec:simulation-model-components}
%%%%%%%%%%%%%%%%%%%%%%%%%%%%%%%%%%%%%%%%%%%%%%%%%%%%%%%%%%%%

A simulation typically includes:

\[
\boxed{
\text{state variables},
}
\]

\[
\boxed{
\text{parameters},
}
\]

\[
\boxed{
\text{inputs},
}
\]

\[
\boxed{
\text{initial conditions},
}
\]

and

\[
\boxed{
\text{outputs}.
}
\]

These categories should be defined explicitly before implementation.

%%%%%%%%%%%%%%%%%%%%%%%%%%%%%%%%%%%%%%%%%%%%%%%%%%%%%%%%%%%%
\subsection{State Variables}
\label{subsec:simulation-state-variables}
%%%%%%%%%%%%%%%%%%%%%%%%%%%%%%%%%%%%%%%%%%%%%%%%%%%%%%%%%%%%

A state variable contains information required to determine future model
behavior.

If

\[
x(t)
\]

is the state, then a deterministic continuous-time model may have form

\[
\boxed{
\frac{dx}{dt}
=
f(x,t;\theta),
}
\]

where

\[
\theta
\]

denotes model parameters.

%%%%%%%%%%%%%%%%%%%%%%%%%%%%%%%%%%%%%%%%%%%%%%%%%%%%%%%%%%%%
\subsection{Parameters}
\label{subsec:simulation-parameters}
%%%%%%%%%%%%%%%%%%%%%%%%%%%%%%%%%%%%%%%%%%%%%%%%%%%%%%%%%%%%

Parameters describe fixed or externally selected model properties.

Examples include:

\[
\boxed{
\text{growth rate},
}
\]

\[
\boxed{
\text{diffusion coefficient},
}
\]

\[
\boxed{
\text{service rate},
}
\]

and

\[
\boxed{
\text{transmission rate}.
}
\]

A parameter may be treated as fixed in one simulation and uncertain in
another.

%%%%%%%%%%%%%%%%%%%%%%%%%%%%%%%%%%%%%%%%%%%%%%%%%%%%%%%%%%%%
\subsection{Inputs and Forcing}
\label{subsec:simulation-inputs}
%%%%%%%%%%%%%%%%%%%%%%%%%%%%%%%%%%%%%%%%%%%%%%%%%%%%%%%%%%%%

An input is an externally prescribed quantity that influences the state.

A model may be written

\[
\boxed{
\dot{x}
=
f(x,u,t;\theta),
}
\]

where

\[
u(t)
\]

is the input or forcing.

Examples include:

\[
\boxed{
\text{external load},
}
\]

\[
\boxed{
\text{control signal},
}
\]

\[
\boxed{
\text{rainfall},
}
\]

or

\[
\boxed{
\text{arrival process}.
}
\]

%%%%%%%%%%%%%%%%%%%%%%%%%%%%%%%%%%%%%%%%%%%%%%%%%%%%%%%%%%%%
\subsection{Outputs}
\label{subsec:simulation-outputs}
%%%%%%%%%%%%%%%%%%%%%%%%%%%%%%%%%%%%%%%%%%%%%%%%%%%%%%%%%%%%

The full state may not be the final quantity of interest.

An output may be

\[
\boxed{
y(t)
=
g(x(t),\theta).
}
\]

Examples include:

\[
\boxed{
\text{maximum temperature},
}
\]

\[
\boxed{
\text{total deaths},
}
\]

\[
\boxed{
\text{waiting time},
}
\]

or

\[
\boxed{
\text{financial loss}.
}
\]

Simulation design should begin by identifying the quantities of interest.

%%%%%%%%%%%%%%%%%%%%%%%%%%%%%%%%%%%%%%%%%%%%%%%%%%%%%%%%%%%%
\subsection{Deterministic Simulation}
\label{subsec:deterministic-simulation}
%%%%%%%%%%%%%%%%%%%%%%%%%%%%%%%%%%%%%%%%%%%%%%%%%%%%%%%%%%%%

A deterministic simulation produces the same mathematical trajectory from
the same:

\[
\boxed{
\text{initial conditions},
}
\]

\[
\boxed{
\text{parameters},
}
\]

and

\[
\boxed{
\text{inputs}.
}
\]

For example,

\[
\dot{x}
=
f(x,t)
\]

with a fixed initial condition defines one deterministic trajectory when
the solution is unique.

%%%%%%%%%%%%%%%%%%%%%%%%%%%%%%%%%%%%%%%%%%%%%%%%%%%%%%%%%%%%
\subsection{Stochastic Simulation}
\label{subsec:stochastic-simulation}
%%%%%%%%%%%%%%%%%%%%%%%%%%%%%%%%%%%%%%%%%%%%%%%%%%%%%%%%%%%%

A stochastic simulation includes random variables or stochastic processes.

The output is therefore random.

One run produces one realization:

\[
\boxed{
X^{(1)}(t).
}
\]

Repeated runs produce

\[
\boxed{
X^{(1)}(t),
X^{(2)}(t),
\ldots,
X^{(N)}(t).
}
\]

Statistical summaries are then used to describe simulation behavior.

%%%%%%%%%%%%%%%%%%%%%%%%%%%%%%%%%%%%%%%%%%%%%%%%%%%%%%%%%%%%
\subsection{Static Simulation}
\label{subsec:static-simulation}
%%%%%%%%%%%%%%%%%%%%%%%%%%%%%%%%%%%%%%%%%%%%%%%%%%%%%%%%%%%%

A static simulation does not explicitly model evolution through time.

For example, if uncertain input

\[
X
\]

produces output

\[
Y=g(X),
\]

then repeated sampling of \(X\) creates a static Monte Carlo simulation.

%%%%%%%%%%%%%%%%%%%%%%%%%%%%%%%%%%%%%%%%%%%%%%%%%%%%%%%%%%%%
\subsection{Dynamic Simulation}
\label{subsec:dynamic-simulation}
%%%%%%%%%%%%%%%%%%%%%%%%%%%%%%%%%%%%%%%%%%%%%%%%%%%%%%%%%%%%

A dynamic simulation models evolution through time.

Examples include:

\[
\boxed{
\text{ODE trajectories},
}
\]

\[
\boxed{
\text{PDE fields},
}
\]

\[
\boxed{
\text{queues},
}
\]

and

\[
\boxed{
\text{agent populations}.
}
\]

The state changes from one time to another.

%%%%%%%%%%%%%%%%%%%%%%%%%%%%%%%%%%%%%%%%%%%%%%%%%%%%%%%%%%%%
\subsection{Continuous-Time Simulation}
\label{subsec:continuous-time-simulation}
%%%%%%%%%%%%%%%%%%%%%%%%%%%%%%%%%%%%%%%%%%%%%%%%%%%%%%%%%%%%

A continuous-time model may have form

\[
\boxed{
\dot{x}
=
f(x,t).
}
\]

The mathematical state is defined continuously in time.

The numerical simulation still evaluates the state only at finitely many
times, using methods such as those of
Section~\ref{sec:numerical-odes}.

%%%%%%%%%%%%%%%%%%%%%%%%%%%%%%%%%%%%%%%%%%%%%%%%%%%%%%%%%%%%
\subsection{Discrete-Time Simulation}
\label{subsec:discrete-time-simulation}
%%%%%%%%%%%%%%%%%%%%%%%%%%%%%%%%%%%%%%%%%%%%%%%%%%%%%%%%%%%%

A discrete-time model evolves through a recurrence:

\[
\boxed{
x_{n+1}
=
F(x_n,n).
}
\]

Examples include:

\[
\boxed{
\text{population maps},
}
\]

\[
\boxed{
\text{difference equations},
}
\]

and

\[
\boxed{
\text{discrete control systems}.
}
\]

No numerical time-discretization is required if the original mathematical
model is itself discrete.

%%%%%%%%%%%%%%%%%%%%%%%%%%%%%%%%%%%%%%%%%%%%%%%%%%%%%%%%%%%%
\subsection{Worked Example: Logistic Map}
\label{subsec:worked-logistic-map-simulation}
%%%%%%%%%%%%%%%%%%%%%%%%%%%%%%%%%%%%%%%%%%%%%%%%%%%%%%%%%%%%

\begin{workedexamplebox}{Discrete Population Simulation}

Consider

\[
x_{n+1}
=
rx_n(1-x_n),
\]

with

\[
r=2,
\qquad
x_0=0.2.
\]

Then

\[
x_1
=
2(0.2)(0.8)
=
0.32.
\]

Next,

\[
x_2
=
2(0.32)(0.68)
=
0.4352.
\]

Then,

\[
x_3
=
2(0.4352)(0.5648)
\approx
0.4916.
\]

\begin{answerbox}
\[
\boxed{
x_1=0.32,
\qquad
x_2=0.4352,
\qquad
x_3\approx0.4916.
}
\]
\end{answerbox}

\end{workedexamplebox}

%%%%%%%%%%%%%%%%%%%%%%%%%%%%%%%%%%%%%%%%%%%%%%%%%%%%%%%%%%%%
\subsection{Time-Driven Simulation}
\label{subsec:time-driven-simulation}
%%%%%%%%%%%%%%%%%%%%%%%%%%%%%%%%%%%%%%%%%%%%%%%%%%%%%%%%%%%%

A time-driven simulation advances using regularly or adaptively selected
time steps:

\[
\boxed{
t_0
\rightarrow
t_1
\rightarrow
t_2
\rightarrow
\cdots.
}
\]

At each step, the state is updated.

Numerical ODE and PDE solvers are standard examples.

%%%%%%%%%%%%%%%%%%%%%%%%%%%%%%%%%%%%%%%%%%%%%%%%%%%%%%%%%%%%
\subsection{Event-Driven Simulation}
\label{subsec:event-driven-simulation}
%%%%%%%%%%%%%%%%%%%%%%%%%%%%%%%%%%%%%%%%%%%%%%%%%%%%%%%%%%%%

In an event-driven simulation, the clock jumps from one event time to the
next.

Suppose event times are

\[
\boxed{
T_1<T_2<T_3<\cdots.
}
\]

The state changes only when an event occurs.

This approach is common in:

\[
\boxed{
\text{queueing},
}
\]

\[
\boxed{
\text{inventory},
}
\]

and

\[
\boxed{
\text{communication networks}.
}
\]

%%%%%%%%%%%%%%%%%%%%%%%%%%%%%%%%%%%%%%%%%%%%%%%%%%%%%%%%%%%%
\subsection{Event Calendar}
\label{subsec:event-calendar}
%%%%%%%%%%%%%%%%%%%%%%%%%%%%%%%%%%%%%%%%%%%%%%%%%%%%%%%%%%%%

A discrete-event simulator maintains a future-event list.

Each event record may contain:

\[
\boxed{
\text{event time},
}
\]

\[
\boxed{
\text{event type},
}
\]

and

\[
\boxed{
\text{associated data}.
}
\]

The next event is removed from the calendar, processed, and any newly
scheduled events are inserted.

%%%%%%%%%%%%%%%%%%%%%%%%%%%%%%%%%%%%%%%%%%%%%%%%%%%%%%%%%%%%
\subsection{Queue Simulation}
\label{subsec:queue-simulation}
%%%%%%%%%%%%%%%%%%%%%%%%%%%%%%%%%%%%%%%%%%%%%%%%%%%%%%%%%%%%

A queueing simulation may include two event types:

\[
\boxed{
\text{arrival}
}
\]

and

\[
\boxed{
\text{departure}.
}
\]

The state may include:

\[
\boxed{
Q(t)
=
\text{number waiting},
}
\]

and the server status.

Events update the state and schedule future events.

%%%%%%%%%%%%%%%%%%%%%%%%%%%%%%%%%%%%%%%%%%%%%%%%%%%%%%%%%%%%
\subsection{Worked Example: Simple Queue Events}
\label{subsec:worked-simple-queue-events}
%%%%%%%%%%%%%%%%%%%%%%%%%%%%%%%%%%%%%%%%%%%%%%%%%%%%%%%%%%%%

\begin{workedexamplebox}{Processing Arrivals and Departures}

Suppose a single-server system is initially empty.

The first arrival occurs at

\[
t=1.
\]

The customer begins service immediately.

Suppose service takes

\[
2
\]

time units.

Then the departure is scheduled at

\[
t=3.
\]

If another customer arrives at

\[
t=2,
\]

the second customer must wait.

\begin{answerbox}
At time \(t=2\),

\[
\boxed{
\text{one customer is in service and one customer is waiting.}
}
\]
\end{answerbox}

\end{workedexamplebox}

%%%%%%%%%%%%%%%%%%%%%%%%%%%%%%%%%%%%%%%%%%%%%%%%%%%%%%%%%%%%
\subsection{Agent-Based Simulation}
\label{subsec:agent-based-simulation}
%%%%%%%%%%%%%%%%%%%%%%%%%%%%%%%%%%%%%%%%%%%%%%%%%%%%%%%%%%%%

Agent-based models simulate individual interacting entities.

An agent may have:

\[
\boxed{
\text{state},
}
\]

\[
\boxed{
\text{rules},
}
\]

\[
\boxed{
\text{location},
}
\]

and

\[
\boxed{
\text{interactions}.
}
\]

System-level behavior emerges from many local interactions.

%%%%%%%%%%%%%%%%%%%%%%%%%%%%%%%%%%%%%%%%%%%%%%%%%%%%%%%%%%%%
\subsection{Examples of Agents}
\label{subsec:simulation-agent-examples}
%%%%%%%%%%%%%%%%%%%%%%%%%%%%%%%%%%%%%%%%%%%%%%%%%%%%%%%%%%%%

Agents may represent:

\[
\boxed{
\text{people},
}
\]

\[
\boxed{
\text{animals},
}
\]

\[
\boxed{
\text{vehicles},
}
\]

\[
\boxed{
\text{firms},
}
\]

or

\[
\boxed{
\text{cells}.
}
\]

Agent-based modeling is useful when individual heterogeneity and local
interaction are important.

%%%%%%%%%%%%%%%%%%%%%%%%%%%%%%%%%%%%%%%%%%%%%%%%%%%%%%%%%%%%
\subsection{Synchronous and Asynchronous Agent Updates}
\label{subsec:agent-update-rules}
%%%%%%%%%%%%%%%%%%%%%%%%%%%%%%%%%%%%%%%%%%%%%%%%%%%%%%%%%%%%

In synchronous updating, every agent computes its next state from the same
current configuration.

Then all states are updated simultaneously.

In asynchronous updating, agents update one at a time or at random times.

The update rule can materially change model behavior.

%%%%%%%%%%%%%%%%%%%%%%%%%%%%%%%%%%%%%%%%%%%%%%%%%%%%%%%%%%%%
\subsection{Compartment Models}
\label{subsec:simulation-compartment-models}
%%%%%%%%%%%%%%%%%%%%%%%%%%%%%%%%%%%%%%%%%%%%%%%%%%%%%%%%%%%%

Compartment models divide a population or quantity among aggregate states.

For example,

\[
\boxed{
S+I+R=N.
}
\]

An epidemic model may use

\[
\boxed{
\frac{dS}{dt}
=
-\beta
\frac{SI}{N},
}
\]

\[
\boxed{
\frac{dI}{dt}
=
\beta
\frac{SI}{N}
-
\gamma I,
}
\]

and

\[
\boxed{
\frac{dR}{dt}
=
\gamma I.
}
\]

The simulation tracks movement among compartments.

%%%%%%%%%%%%%%%%%%%%%%%%%%%%%%%%%%%%%%%%%%%%%%%%%%%%%%%%%%%%
\subsection{Conservation in Compartment Models}
\label{subsec:compartment-conservation}
%%%%%%%%%%%%%%%%%%%%%%%%%%%%%%%%%%%%%%%%%%%%%%%%%%%%%%%%%%%%

For the SIR equations,

\[
\frac{d}{dt}
(S+I+R)
=
0.
\]

Therefore,

\[
\boxed{
S(t)+I(t)+R(t)=N.
}
\]

A numerical simulation should preserve this total exactly or to within
small numerical error, depending on the numerical method.

This provides a useful verification diagnostic.

%%%%%%%%%%%%%%%%%%%%%%%%%%%%%%%%%%%%%%%%%%%%%%%%%%%%%%%%%%%%
\subsection{Network Simulation}
\label{subsec:network-simulation}
%%%%%%%%%%%%%%%%%%%%%%%%%%%%%%%%%%%%%%%%%%%%%%%%%%%%%%%%%%%%

Many systems are naturally represented by a graph

\[
\boxed{
G=(V,E).
}
\]

Vertices represent entities and edges represent interactions.

A network simulation may model:

\[
\boxed{
\text{disease transmission},
}
\]

\[
\boxed{
\text{traffic},
}
\]

\[
\boxed{
\text{power flow},
}
\]

or

\[
\boxed{
\text{information spread}.
}
\]

%%%%%%%%%%%%%%%%%%%%%%%%%%%%%%%%%%%%%%%%%%%%%%%%%%%%%%%%%%%%
\subsection{State on a Network}
\label{subsec:state-on-network}
%%%%%%%%%%%%%%%%%%%%%%%%%%%%%%%%%%%%%%%%%%%%%%%%%%%%%%%%%%%%

Suppose each node \(i\) has state

\[
x_i(t).
\]

A coupled system may have form

\[
\boxed{
\dot{x}_i
=
f_i(x_i)
+
\sum_{j}
A_{ij}
g_{ij}(x_i,x_j),
}
\]

where

\[
A
\]

is an adjacency matrix.

The network structure controls which states interact.

%%%%%%%%%%%%%%%%%%%%%%%%%%%%%%%%%%%%%%%%%%%%%%%%%%%%%%%%%%%%
\subsection{Cellular Automata}
\label{subsec:cellular-automata-simulation}
%%%%%%%%%%%%%%%%%%%%%%%%%%%%%%%%%%%%%%%%%%%%%%%%%%%%%%%%%%%%

A cellular automaton consists of:

\[
\boxed{
\text{a lattice},
}
\]

\[
\boxed{
\text{discrete states},
}
\]

and

\[
\boxed{
\text{local update rules}.
}
\]

At each discrete time, the state of each cell is updated based on a
neighborhood.

Simple local rules can generate complex global patterns.

%%%%%%%%%%%%%%%%%%%%%%%%%%%%%%%%%%%%%%%%%%%%%%%%%%%%%%%%%%%%
\subsection{State-Transition Simulation}
\label{subsec:state-transition-simulation}
%%%%%%%%%%%%%%%%%%%%%%%%%%%%%%%%%%%%%%%%%%%%%%%%%%%%%%%%%%%%

A general discrete stochastic model may use transition probabilities

\[
\boxed{
P_{ij}
=
P(X_{n+1}=j\mid X_n=i).
}
\]

Simulation then draws the next state from the row of \(P\) corresponding
to the current state.

This is a Markov-chain simulation.

%%%%%%%%%%%%%%%%%%%%%%%%%%%%%%%%%%%%%%%%%%%%%%%%%%%%%%%%%%%%
\subsection{Simulation Horizon}
\label{subsec:simulation-horizon}
%%%%%%%%%%%%%%%%%%%%%%%%%%%%%%%%%%%%%%%%%%%%%%%%%%%%%%%%%%%%

A simulation must define a stopping horizon.

Examples include:

\[
\boxed{
0\leq t\leq T,
}
\]

\[
\boxed{
n=0,\ldots,N,
}
\]

or

\[
\boxed{
\text{until a stopping event occurs}.
}
\]

The horizon should be long enough to capture the behavior relevant to the
scientific question.

%%%%%%%%%%%%%%%%%%%%%%%%%%%%%%%%%%%%%%%%%%%%%%%%%%%%%%%%%%%%
\subsection{Transient Behavior}
\label{subsec:simulation-transient-behavior}
%%%%%%%%%%%%%%%%%%%%%%%%%%%%%%%%%%%%%%%%%%%%%%%%%%%%%%%%%%%%

Early simulation behavior can depend strongly on initial conditions.

This initial period is called the transient.

For example, a queue initialized empty may initially have artificially
short waiting times.

Steady-state analysis may therefore discard or separately analyze an
initial warm-up period.

%%%%%%%%%%%%%%%%%%%%%%%%%%%%%%%%%%%%%%%%%%%%%%%%%%%%%%%%%%%%
\subsection{Steady-State Simulation}
\label{subsec:steady-state-simulation}
%%%%%%%%%%%%%%%%%%%%%%%%%%%%%%%%%%%%%%%%%%%%%%%%%%%%%%%%%%%%

A stochastic system may approach a stationary distribution.

If

\[
X(t)
\]

reaches statistical equilibrium, long-run quantities may be estimated from
a sufficiently long simulation.

Examples include:

\[
\boxed{
\text{mean queue length},
}
\]

\[
\boxed{
\text{server utilization},
}
\]

and

\[
\boxed{
\text{long-run inventory level}.
}
\]

%%%%%%%%%%%%%%%%%%%%%%%%%%%%%%%%%%%%%%%%%%%%%%%%%%%%%%%%%%%%
\subsection{Warm-Up Period}
\label{subsec:simulation-warmup-period}
%%%%%%%%%%%%%%%%%%%%%%%%%%%%%%%%%%%%%%%%%%%%%%%%%%%%%%%%%%%%

Suppose observations during

\[
0\leq t<T_w
\]

are strongly influenced by artificial initial conditions.

One may discard this period and estimate steady-state quantities using

\[
\boxed{
t\geq T_w.
}
\]

The warm-up length should be justified rather than selected arbitrarily.

%%%%%%%%%%%%%%%%%%%%%%%%%%%%%%%%%%%%%%%%%%%%%%%%%%%%%%%%%%%%
\subsection{Finite-Horizon Simulation}
\label{subsec:finite-horizon-simulation}
%%%%%%%%%%%%%%%%%%%%%%%%%%%%%%%%%%%%%%%%%%%%%%%%%%%%%%%%%%%%

Some systems are inherently finite-horizon.

Examples include:

\[
\boxed{
\text{one epidemic season},
}
\]

\[
\boxed{
\text{one flight operation},
}
\]

\[
\boxed{
\text{a manufacturing shift},
}
\]

or

\[
\boxed{
\text{a financial option maturity period}.
}
\]

In these cases, transient behavior may itself be the quantity of interest.

%%%%%%%%%%%%%%%%%%%%%%%%%%%%%%%%%%%%%%%%%%%%%%%%%%%%%%%%%%%%
\subsection{Time-Step Selection}
\label{subsec:simulation-time-step-selection}
%%%%%%%%%%%%%%%%%%%%%%%%%%%%%%%%%%%%%%%%%%%%%%%%%%%%%%%%%%%%

For time-driven numerical simulations, the time step must balance:

\[
\boxed{
\text{accuracy},
}
\]

\[
\boxed{
\text{stability},
}
\]

and

\[
\boxed{
\text{cost}.
}
\]

A time step that is too large may miss important dynamics.

A time step that is excessively small may create unnecessary computation.

%%%%%%%%%%%%%%%%%%%%%%%%%%%%%%%%%%%%%%%%%%%%%%%%%%%%%%%%%%%%
\subsection{Physical and Numerical Time Scales}
\label{subsec:physical-numerical-time-scales}
%%%%%%%%%%%%%%%%%%%%%%%%%%%%%%%%%%%%%%%%%%%%%%%%%%%%%%%%%%%%

Suppose the model contains characteristic time scales

\[
\boxed{
\tau_1,
\tau_2,\ldots.
}
\]

The numerical method must resolve the time scales relevant to the desired
output.

Fast time scales may also create stiffness, requiring special integration
methods.

%%%%%%%%%%%%%%%%%%%%%%%%%%%%%%%%%%%%%%%%%%%%%%%%%%%%%%%%%%%%
\subsection{Physical Instability Versus Numerical Instability}
\label{subsec:physical-vs-numerical-instability}
%%%%%%%%%%%%%%%%%%%%%%%%%%%%%%%%%%%%%%%%%%%%%%%%%%%%%%%%%%%%

A physical model may genuinely exhibit exponential growth.

A numerical method may also produce artificial growth even when the exact
solution is stable.

These phenomena must be distinguished.

A convergence or stability study can help determine whether observed
growth is physical or numerical.

%%%%%%%%%%%%%%%%%%%%%%%%%%%%%%%%%%%%%%%%%%%%%%%%%%%%%%%%%%%%
\subsection{Initial Conditions}
\label{subsec:simulation-initial-conditions}
%%%%%%%%%%%%%%%%%%%%%%%%%%%%%%%%%%%%%%%%%%%%%%%%%%%%%%%%%%%%

Initial conditions can strongly influence simulation outcomes.

A simulation study may consider:

\[
\boxed{
x(0)=x_0,
}
\]

or a distribution

\[
\boxed{
X(0)\sim p_0.
}
\]

When multiple plausible initial states exist, sensitivity to initial
conditions should be investigated.

%%%%%%%%%%%%%%%%%%%%%%%%%%%%%%%%%%%%%%%%%%%%%%%%%%%%%%%%%%%%
\subsection{Chaotic Sensitivity}
\label{subsec:chaotic-simulation-sensitivity}
%%%%%%%%%%%%%%%%%%%%%%%%%%%%%%%%%%%%%%%%%%%%%%%%%%%%%%%%%%%%

In chaotic systems, nearby initial conditions can separate rapidly.

If

\[
\|x_0-\widetilde{x}_0\|
\]

is small, it is possible that

\[
\boxed{
\|x(t)-\widetilde{x}(t)\|
}
\]

becomes large at later times.

Long-term trajectory prediction may therefore be impossible even when the
model is deterministic.

Statistical or qualitative properties may remain meaningful.

%%%%%%%%%%%%%%%%%%%%%%%%%%%%%%%%%%%%%%%%%%%%%%%%%%%%%%%%%%%%
\subsection{Conservation Laws as Diagnostics}
\label{subsec:simulation-conservation-diagnostics}
%%%%%%%%%%%%%%%%%%%%%%%%%%%%%%%%%%%%%%%%%%%%%%%%%%%%%%%%%%%%

Suppose the model conserves

\[
I(x).
\]

Then

\[
\boxed{
I(x(t))
=
I(x(0)).
}
\]

Monitor

\[
\boxed{
\Delta I(t)
=
I(x(t))-I(x(0)).
}
\]

Unexpected drift can reveal numerical error or implementation problems.

%%%%%%%%%%%%%%%%%%%%%%%%%%%%%%%%%%%%%%%%%%%%%%%%%%%%%%%%%%%%
\subsection{Positivity Diagnostics}
\label{subsec:simulation-positivity}
%%%%%%%%%%%%%%%%%%%%%%%%%%%%%%%%%%%%%%%%%%%%%%%%%%%%%%%%%%%%

Many model quantities must remain nonnegative:

\[
\boxed{
x_i(t)\geq0.
}
\]

Examples include:

\[
\boxed{
\text{population},
}
\]

\[
\boxed{
\text{mass},
}
\]

\[
\boxed{
\text{concentration},
}
\]

and

\[
\boxed{
\text{probability}.
}
\]

Negative computed values may indicate inappropriate numerical methods or
step sizes.

%%%%%%%%%%%%%%%%%%%%%%%%%%%%%%%%%%%%%%%%%%%%%%%%%%%%%%%%%%%%
\subsection{Boundedness Diagnostics}
\label{subsec:simulation-boundedness}
%%%%%%%%%%%%%%%%%%%%%%%%%%%%%%%%%%%%%%%%%%%%%%%%%%%%%%%%%%%%

A model may guarantee

\[
\boxed{
a\leq x(t)\leq b.
}
\]

If a simulation violates these known bounds, possible causes include:

\[
\boxed{
\text{numerical instability},
}
\]

\[
\boxed{
\text{coding errors},
}
\]

or

\[
\boxed{
\text{incorrect parameterization}.
}
\]

Known mathematical bounds are valuable computational checks.

%%%%%%%%%%%%%%%%%%%%%%%%%%%%%%%%%%%%%%%%%%%%%%%%%%%%%%%%%%%%
\subsection{Stochastic Replications}
\label{subsec:stochastic-replications}
%%%%%%%%%%%%%%%%%%%%%%%%%%%%%%%%%%%%%%%%%%%%%%%%%%%%%%%%%%%%

A single stochastic trajectory cannot characterize the full behavior of a
random system.

Generate independent replications

\[
\boxed{
Y_1,\ldots,Y_N.
}
\]

Estimate the mean by

\[
\boxed{
\overline{Y}
=
\frac1N
\sum_{i=1}^{N}
Y_i.
}
\]

Estimate the variance by

\[
\boxed{
s^2
=
\frac1{N-1}
\sum_{i=1}^{N}
(Y_i-\overline{Y})^2.
}
\]

%%%%%%%%%%%%%%%%%%%%%%%%%%%%%%%%%%%%%%%%%%%%%%%%%%%%%%%%%%%%
\subsection{Ensemble Mean}
\label{subsec:simulation-ensemble-mean}
%%%%%%%%%%%%%%%%%%%%%%%%%%%%%%%%%%%%%%%%%%%%%%%%%%%%%%%%%%%%

Suppose each replication produces a trajectory

\[
X^{(r)}(t).
\]

The ensemble mean is

\[
\boxed{
\overline{X}(t)
=
\frac1R
\sum_{r=1}^{R}
X^{(r)}(t).
}
\]

This estimates

\[
\boxed{
\mathbb{E}[X(t)].
}
\]

Pointwise quantiles can also be estimated across the ensemble.

%%%%%%%%%%%%%%%%%%%%%%%%%%%%%%%%%%%%%%%%%%%%%%%%%%%%%%%%%%%%
\subsection{Simulation Confidence Bands}
\label{subsec:simulation-confidence-bands}
%%%%%%%%%%%%%%%%%%%%%%%%%%%%%%%%%%%%%%%%%%%%%%%%%%%%%%%%%%%%

At each time \(t\), repeated simulations can estimate uncertainty in a
summary such as the mean.

For independent runs, an approximate pointwise interval is

\[
\boxed{
\overline{X}(t)
\pm
z_{1-\alpha/2}
\frac{
s(t)
}{
\sqrt{R}
}.
}
\]

These are pointwise intervals, not automatically simultaneous confidence
bands over the entire trajectory.

%%%%%%%%%%%%%%%%%%%%%%%%%%%%%%%%%%%%%%%%%%%%%%%%%%%%%%%%%%%%
\subsection{Aleatoric Uncertainty}
\label{subsec:simulation-aleatoric-uncertainty}
%%%%%%%%%%%%%%%%%%%%%%%%%%%%%%%%%%%%%%%%%%%%%%%%%%%%%%%%%%%%

Aleatoric uncertainty represents intrinsic randomness in the modeled
system.

Examples include:

\[
\boxed{
\text{random arrivals},
}
\]

\[
\boxed{
\text{thermal fluctuations},
}
\]

and

\[
\boxed{
\text{random failure times}.
}
\]

Increasing the amount of data may characterize this randomness more
precisely, but does not eliminate it from the system.

%%%%%%%%%%%%%%%%%%%%%%%%%%%%%%%%%%%%%%%%%%%%%%%%%%%%%%%%%%%%
\subsection{Epistemic or Parametric Uncertainty}
\label{subsec:simulation-parametric-uncertainty}
%%%%%%%%%%%%%%%%%%%%%%%%%%%%%%%%%%%%%%%%%%%%%%%%%%%%%%%%%%%%

Parameter uncertainty arises because quantities such as

\[
\theta
\]

are not known exactly.

One may represent

\[
\boxed{
\Theta
\sim p(\theta)
}
\]

and propagate parameter uncertainty through the simulation.

This is different from intrinsic randomness in the system dynamics.

%%%%%%%%%%%%%%%%%%%%%%%%%%%%%%%%%%%%%%%%%%%%%%%%%%%%%%%%%%%%
\subsection{Uncertainty Propagation}
\label{subsec:simulation-uncertainty-propagation}
%%%%%%%%%%%%%%%%%%%%%%%%%%%%%%%%%%%%%%%%%%%%%%%%%%%%%%%%%%%%

Let

\[
Y
=
M(\Theta).
\]

Sample

\[
\Theta_1,\ldots,\Theta_N
\]

from the uncertainty model.

Then compute

\[
\boxed{
Y_i
=
M(\Theta_i).
}
\]

The resulting outputs approximate the distribution of \(Y\).

%%%%%%%%%%%%%%%%%%%%%%%%%%%%%%%%%%%%%%%%%%%%%%%%%%%%%%%%%%%%
\subsection{Output Distribution}
\label{subsec:simulation-output-distribution}
%%%%%%%%%%%%%%%%%%%%%%%%%%%%%%%%%%%%%%%%%%%%%%%%%%%%%%%%%%%%

Simulation outputs can be summarized by:

\[
\boxed{
\mathbb{E}[Y],
}
\]

\[
\boxed{
\operatorname{Var}(Y),
}
\]

\[
\boxed{
P(Y>c),
}
\]

and quantiles

\[
\boxed{
q_\alpha.
}
\]

Reporting only the average may hide important tail behavior or
multimodality.

%%%%%%%%%%%%%%%%%%%%%%%%%%%%%%%%%%%%%%%%%%%%%%%%%%%%%%%%%%%%
\subsection{Scenario Analysis}
\label{subsec:simulation-scenario-analysis}
%%%%%%%%%%%%%%%%%%%%%%%%%%%%%%%%%%%%%%%%%%%%%%%%%%%%%%%%%%%%

Scenario analysis compares selected parameter or policy settings.

For example:

\[
\boxed{
\text{baseline},
}
\]

\[
\boxed{
\text{moderate intervention},
}
\]

and

\[
\boxed{
\text{strong intervention}.
}
\]

Each scenario should be defined explicitly and simulated using a consistent
comparison framework.

%%%%%%%%%%%%%%%%%%%%%%%%%%%%%%%%%%%%%%%%%%%%%%%%%%%%%%%%%%%%
\subsection{Controlled Scenario Comparison}
\label{subsec:controlled-scenario-comparison}
%%%%%%%%%%%%%%%%%%%%%%%%%%%%%%%%%%%%%%%%%%%%%%%%%%%%%%%%%%%%

When comparing scenarios, keep unrelated assumptions fixed.

Then observed changes can be attributed more clearly to the changed
scenario variable.

In stochastic simulation, common random numbers may further reduce noise
in pairwise comparisons.

%%%%%%%%%%%%%%%%%%%%%%%%%%%%%%%%%%%%%%%%%%%%%%%%%%%%%%%%%%%%
\subsection{Parameter Sweeps}
\label{subsec:simulation-parameter-sweeps}
%%%%%%%%%%%%%%%%%%%%%%%%%%%%%%%%%%%%%%%%%%%%%%%%%%%%%%%%%%%%

Suppose the model depends on parameter pair

\[
(\theta_1,\theta_2).
\]

Evaluate the simulation over a grid:

\[
\boxed{
(\theta_{1,i},\theta_{2,j}).
}
\]

For each pair, compute quantity of interest

\[
\boxed{
Q_{ij}.
}
\]

The resulting table or surface reveals broad parameter dependence.

%%%%%%%%%%%%%%%%%%%%%%%%%%%%%%%%%%%%%%%%%%%%%%%%%%%%%%%%%%%%
\subsection{One-at-a-Time Sensitivity}
\label{subsec:one-at-a-time-sensitivity}
%%%%%%%%%%%%%%%%%%%%%%%%%%%%%%%%%%%%%%%%%%%%%%%%%%%%%%%%%%%%

A simple sensitivity approach changes one parameter while holding all
others fixed.

For parameter \(\theta_i\),

\[
\boxed{
S_i
\approx
\frac{
Q(\theta_i+h)-Q(\theta_i-h)
}{
2h
}.
}
\]

This estimates a local derivative.

It does not capture interactions between several uncertain parameters.

%%%%%%%%%%%%%%%%%%%%%%%%%%%%%%%%%%%%%%%%%%%%%%%%%%%%%%%%%%%%
\subsection{Normalized Sensitivity}
\label{subsec:normalized-sensitivity}
%%%%%%%%%%%%%%%%%%%%%%%%%%%%%%%%%%%%%%%%%%%%%%%%%%%%%%%%%%%%

A dimensionless local sensitivity may be written

\[
\boxed{
S_i^{\mathrm{rel}}
=
\frac{
\theta_i
}{
Q
}
\frac{
\partial Q
}{
\partial\theta_i
},
}
\]

provided the quantities are nonzero.

This measures approximate percentage output change per percentage parameter
change.

%%%%%%%%%%%%%%%%%%%%%%%%%%%%%%%%%%%%%%%%%%%%%%%%%%%%%%%%%%%%
\subsection{Global Sensitivity}
\label{subsec:simulation-global-sensitivity}
%%%%%%%%%%%%%%%%%%%%%%%%%%%%%%%%%%%%%%%%%%%%%%%%%%%%%%%%%%%%

Global sensitivity analysis varies parameters over their entire plausible
ranges.

It can capture:

\[
\boxed{
\text{nonlinearity},
}
\]

\[
\boxed{
\text{interactions},
}
\]

and

\[
\boxed{
\text{nonlocal effects}.
}
\]

Monte Carlo sampling is commonly used for global sensitivity studies.

%%%%%%%%%%%%%%%%%%%%%%%%%%%%%%%%%%%%%%%%%%%%%%%%%%%%%%%%%%%%
\subsection{Variance-Based Sensitivity Preview}
\label{subsec:variance-based-sensitivity-preview}
%%%%%%%%%%%%%%%%%%%%%%%%%%%%%%%%%%%%%%%%%%%%%%%%%%%%%%%%%%%%

Suppose

\[
Y=M(\Theta_1,\ldots,\Theta_d).
\]

A variance-based method decomposes

\[
\operatorname{Var}(Y)
\]

into contributions associated with individual parameters and interactions.

A parameter with a large contribution strongly influences output
uncertainty.

%%%%%%%%%%%%%%%%%%%%%%%%%%%%%%%%%%%%%%%%%%%%%%%%%%%%%%%%%%%%
\subsection{Calibration}
\label{subsec:simulation-calibration}
%%%%%%%%%%%%%%%%%%%%%%%%%%%%%%%%%%%%%%%%%%%%%%%%%%%%%%%%%%%%

Calibration chooses parameter values so simulation outputs agree with
observations.

Suppose observed data are

\[
y_1,\ldots,y_m,
\]

and simulated predictions are

\[
\widehat{y}_i(\theta).
\]

A least-squares calibration problem is

\[
\boxed{
\min_\theta
\sum_{i=1}^{m}
\left[
y_i-\widehat{y}_i(\theta)
\right]^2.
}
\]

%%%%%%%%%%%%%%%%%%%%%%%%%%%%%%%%%%%%%%%%%%%%%%%%%%%%%%%%%%%%
\subsection{Weighted Calibration}
\label{subsec:weighted-calibration}
%%%%%%%%%%%%%%%%%%%%%%%%%%%%%%%%%%%%%%%%%%%%%%%%%%%%%%%%%%%%

If observations have different uncertainties, use weights:

\[
\boxed{
J(\theta)
=
\sum_{i=1}^{m}
w_i
\left[
y_i-\widehat{y}_i(\theta)
\right]^2.
}
\]

A common statistical choice is

\[
\boxed{
w_i
=
\frac1{\sigma_i^2}
}
\]

when observational error variances \(\sigma_i^2\) are known or estimated.

%%%%%%%%%%%%%%%%%%%%%%%%%%%%%%%%%%%%%%%%%%%%%%%%%%%%%%%%%%%%
\subsection{Worked Example: Parameter Calibration}
\label{subsec:worked-parameter-calibration}
%%%%%%%%%%%%%%%%%%%%%%%%%%%%%%%%%%%%%%%%%%%%%%%%%%%%%%%%%%%%

\begin{workedexamplebox}{Estimating an Exponential Growth Rate}

Suppose

\[
P(t)
=
P_0e^{rt},
\]

with known

\[
P_0=100.
\]

An observation at

\[
t=2
\]

gives

\[
P(2)=200.
\]

Then

\[
200
=
100e^{2r}.
\]

Thus,

\[
2=e^{2r}.
\]

Taking logarithms,

\[
\ln2=2r.
\]

Therefore,

\begin{answerbox}
\[
\boxed{
r
=
\frac{\ln2}{2}.
}
\]
\end{answerbox}

For more complicated models, the same calibration idea may require a
numerical optimizer.

\end{workedexamplebox}

%%%%%%%%%%%%%%%%%%%%%%%%%%%%%%%%%%%%%%%%%%%%%%%%%%%%%%%%%%%%
\subsection{Simulation-Based Parameter Estimation}
\label{subsec:simulation-based-parameter-estimation}
%%%%%%%%%%%%%%%%%%%%%%%%%%%%%%%%%%%%%%%%%%%%%%%%%%%%%%%%%%%%

If the model has no closed-form solution, computing

\[
J(\theta)
\]

may itself require running a simulation.

Thus calibration becomes

\[
\boxed{
\theta
\rightarrow
\text{simulation}
\rightarrow
\widehat{y}(\theta)
\rightarrow
J(\theta)
\rightarrow
\text{optimizer}.
}
\]

This nested structure can be computationally expensive.

%%%%%%%%%%%%%%%%%%%%%%%%%%%%%%%%%%%%%%%%%%%%%%%%%%%%%%%%%%%%
\subsection{Inverse Problems}
\label{subsec:simulation-inverse-problems}
%%%%%%%%%%%%%%%%%%%%%%%%%%%%%%%%%%%%%%%%%%%%%%%%%%%%%%%%%%%%

A forward problem computes outputs from parameters:

\[
\boxed{
\theta
\longrightarrow
y.
}
\]

An inverse problem attempts to infer parameters from observations:

\[
\boxed{
y
\longrightarrow
\theta.
}
\]

Simulation-based calibration is therefore a form of inverse problem.

%%%%%%%%%%%%%%%%%%%%%%%%%%%%%%%%%%%%%%%%%%%%%%%%%%%%%%%%%%%%
\subsection{Identifiability}
\label{subsec:simulation-identifiability}
%%%%%%%%%%%%%%%%%%%%%%%%%%%%%%%%%%%%%%%%%%%%%%%%%%%%%%%%%%%%

A parameter is identifiable if distinct parameter values produce
distinguishable model outputs in principle.

If

\[
\boxed{
M(\theta_1)
=
M(\theta_2)
}
\]

for

\[
\theta_1\neq\theta_2,
\]

then the observations may not uniquely determine the parameter.

%%%%%%%%%%%%%%%%%%%%%%%%%%%%%%%%%%%%%%%%%%%%%%%%%%%%%%%%%%%%
\subsection{Structural Identifiability}
\label{subsec:structural-identifiability}
%%%%%%%%%%%%%%%%%%%%%%%%%%%%%%%%%%%%%%%%%%%%%%%%%%%%%%%%%%%%

Structural identifiability concerns whether parameters can be uniquely
determined from ideal noise-free observations with perfect experimental
information.

It is a mathematical property of the model structure.

A model may fail structurally even before measurement noise is considered.

%%%%%%%%%%%%%%%%%%%%%%%%%%%%%%%%%%%%%%%%%%%%%%%%%%%%%%%%%%%%
\subsection{Practical Identifiability}
\label{subsec:practical-identifiability}
%%%%%%%%%%%%%%%%%%%%%%%%%%%%%%%%%%%%%%%%%%%%%%%%%%%%%%%%%%%%

A structurally identifiable parameter may still be difficult to estimate
from finite noisy data.

This is practical nonidentifiability.

Signs include:

\[
\boxed{
\text{large confidence intervals},
}
\]

\[
\boxed{
\text{flat objective functions},
}
\]

or

\[
\boxed{
\text{strong parameter correlation}.
}
\]

%%%%%%%%%%%%%%%%%%%%%%%%%%%%%%%%%%%%%%%%%%%%%%%%%%%%%%%%%%%%
\subsection{Calibration Versus Validation}
\label{subsec:calibration-versus-validation}
%%%%%%%%%%%%%%%%%%%%%%%%%%%%%%%%%%%%%%%%%%%%%%%%%%%%%%%%%%%%

Calibration chooses model parameters using data.

Validation evaluates whether the calibrated model predicts independent
observations adequately.

Thus using the same data for both tasks can overstate model performance.

Conceptually,

\[
\boxed{
\text{calibration}
=
\text{fit},
}
\]

while

\[
\boxed{
\text{validation}
=
\text{independent test}.
}
\]

%%%%%%%%%%%%%%%%%%%%%%%%%%%%%%%%%%%%%%%%%%%%%%%%%%%%%%%%%%%%
\subsection{Verification Versus Validation}
\label{subsec:simulation-verification-validation}
%%%%%%%%%%%%%%%%%%%%%%%%%%%%%%%%%%%%%%%%%%%%%%%%%%%%%%%%%%%%

Verification asks:

\[
\boxed{
\text{Did we solve the mathematical model correctly?}
}
\]

Validation asks:

\[
\boxed{
\text{Is the mathematical model adequate for the real system?}
}
\]

These are distinct questions.

A perfectly verified simulation can still be based on a poor model.

%%%%%%%%%%%%%%%%%%%%%%%%%%%%%%%%%%%%%%%%%%%%%%%%%%%%%%%%%%%%
\subsection{Code Verification}
\label{subsec:simulation-code-verification}
%%%%%%%%%%%%%%%%%%%%%%%%%%%%%%%%%%%%%%%%%%%%%%%%%%%%%%%%%%%%

Code verification checks whether the software implements the intended
numerical method correctly.

Methods include:

\[
\boxed{
\text{unit tests},
}
\]

\[
\boxed{
\text{benchmark problems},
}
\]

\[
\boxed{
\text{manufactured solutions},
}
\]

and

\[
\boxed{
\text{convergence studies}.
}
\]

%%%%%%%%%%%%%%%%%%%%%%%%%%%%%%%%%%%%%%%%%%%%%%%%%%%%%%%%%%%%
\subsection{Solution Verification}
\label{subsec:simulation-solution-verification}
%%%%%%%%%%%%%%%%%%%%%%%%%%%%%%%%%%%%%%%%%%%%%%%%%%%%%%%%%%%%

Solution verification estimates numerical uncertainty in one particular
simulation result.

It asks how errors such as

\[
\boxed{
\text{time-step error},
}
\]

\[
\boxed{
\text{mesh error},
}
\]

and

\[
\boxed{
\text{iterative solver error}
}
\]

affect the reported solution.

%%%%%%%%%%%%%%%%%%%%%%%%%%%%%%%%%%%%%%%%%%%%%%%%%%%%%%%%%%%%
\subsection{Benchmark Problems}
\label{subsec:simulation-benchmark-problems}
%%%%%%%%%%%%%%%%%%%%%%%%%%%%%%%%%%%%%%%%%%%%%%%%%%%%%%%%%%%%

A benchmark problem has a trusted analytical or reference solution.

For example, test a numerical ODE solver on

\[
\boxed{
y'=y,
\qquad
y(0)=1,
}
\]

whose exact solution is

\[
\boxed{
y=e^t.
}
\]

Agreement with benchmark behavior provides evidence of implementation
correctness.

%%%%%%%%%%%%%%%%%%%%%%%%%%%%%%%%%%%%%%%%%%%%%%%%%%%%%%%%%%%%
\subsection{Convergence Studies}
\label{subsec:simulation-convergence-studies}
%%%%%%%%%%%%%%%%%%%%%%%%%%%%%%%%%%%%%%%%%%%%%%%%%%%%%%%%%%%%

Suppose a method has expected order \(p\).

Run simulations with

\[
h,
\qquad
\frac h2,
\qquad
\frac h4.
\]

If

\[
E(h)
\approx
Ch^p,
\]

then

\[
\boxed{
\frac{
E(h)
}{
E(h/2)
}
\approx
2^p.
}
\]

Observed convergence provides a strong implementation check.

%%%%%%%%%%%%%%%%%%%%%%%%%%%%%%%%%%%%%%%%%%%%%%%%%%%%%%%%%%%%
\subsection{Synthetic Data}
\label{subsec:simulation-synthetic-data}
%%%%%%%%%%%%%%%%%%%%%%%%%%%%%%%%%%%%%%%%%%%%%%%%%%%%%%%%%%%%

Synthetic data are generated from a known model with known parameters.

For example,

\[
\boxed{
y_i
=
M(t_i;\theta_{\mathrm{true}})
+
\varepsilon_i.
}
\]

A calibration method can then be tested to determine whether it recovers

\[
\theta_{\mathrm{true}}.
\]

This is useful for validating estimation procedures.

%%%%%%%%%%%%%%%%%%%%%%%%%%%%%%%%%%%%%%%%%%%%%%%%%%%%%%%%%%%%
\subsection{Model Validation}
\label{subsec:simulation-model-validation}
%%%%%%%%%%%%%%%%%%%%%%%%%%%%%%%%%%%%%%%%%%%%%%%%%%%%%%%%%%%%

Model validation compares predictions against empirical observations.

Useful comparisons may include:

\[
\boxed{
\text{time-series trajectories},
}
\]

\[
\boxed{
\text{means},
}
\]

\[
\boxed{
\text{variances},
}
\]

\[
\boxed{
\text{distributions},
}
\]

or

\[
\boxed{
\text{event frequencies}.
}
\]

The validation metric should match the intended use of the model.

%%%%%%%%%%%%%%%%%%%%%%%%%%%%%%%%%%%%%%%%%%%%%%%%%%%%%%%%%%%%
\subsection{Residual Analysis}
\label{subsec:simulation-residual-analysis}
%%%%%%%%%%%%%%%%%%%%%%%%%%%%%%%%%%%%%%%%%%%%%%%%%%%%%%%%%%%%

For observations \(y_i\) and predictions \(\widehat{y}_i\), define residuals

\[
\boxed{
r_i
=
y_i-\widehat{y}_i.
}
\]

Residuals can reveal:

\[
\boxed{
\text{bias},
}
\]

\[
\boxed{
\text{time dependence},
}
\]

\[
\boxed{
\text{heteroskedasticity},
}
\]

or

\[
\boxed{
\text{missing model structure}.
}
\]

A small aggregate error can still hide systematic residual patterns.

%%%%%%%%%%%%%%%%%%%%%%%%%%%%%%%%%%%%%%%%%%%%%%%%%%%%%%%%%%%%
\subsection{Simulation Error Decomposition}
\label{subsec:simulation-error-decomposition}
%%%%%%%%%%%%%%%%%%%%%%%%%%%%%%%%%%%%%%%%%%%%%%%%%%%%%%%%%%%%

A discrepancy between simulation and reality may arise from:

\[
\boxed{
\text{measurement error},
}
\]

\[
\boxed{
\text{parameter error},
}
\]

\[
\boxed{
\text{model-form error},
}
\]

\[
\boxed{
\text{discretization error},
}
\]

and

\[
\boxed{
\text{software error}.
}
\]

These sources should not be conflated.

%%%%%%%%%%%%%%%%%%%%%%%%%%%%%%%%%%%%%%%%%%%%%%%%%%%%%%%%%%%%
\subsection{Model-Form Error}
\label{subsec:model-form-error}
%%%%%%%%%%%%%%%%%%%%%%%%%%%%%%%%%%%%%%%%%%%%%%%%%%%%%%%%%%%%

Model-form error occurs when the governing equations omit or simplify
important physical or behavioral mechanisms.

No amount of numerical refinement can remove this error.

Thus

\[
\boxed{
h\to0
}
\]

does not imply that the simulation approaches reality.

It only implies that the numerical solution approaches the chosen
mathematical model.

%%%%%%%%%%%%%%%%%%%%%%%%%%%%%%%%%%%%%%%%%%%%%%%%%%%%%%%%%%%%
\subsection{Uncertainty Quantification}
\label{subsec:simulation-uq}
%%%%%%%%%%%%%%%%%%%%%%%%%%%%%%%%%%%%%%%%%%%%%%%%%%%%%%%%%%%%

Uncertainty quantification studies how uncertain inputs influence
simulation outputs.

Suppose

\[
\Theta
\]

is uncertain and

\[
Y=M(\Theta).
\]

The goal is to characterize the distribution of \(Y\).

This may include:

\[
\boxed{
\mathbb{E}[Y],
}
\]

\[
\boxed{
\operatorname{Var}(Y),
}
\]

\[
\boxed{
P(Y>c),
}
\]

and quantiles.

%%%%%%%%%%%%%%%%%%%%%%%%%%%%%%%%%%%%%%%%%%%%%%%%%%%%%%%%%%%%
\subsection{Confidence and Prediction Intervals}
\label{subsec:simulation-confidence-prediction}
%%%%%%%%%%%%%%%%%%%%%%%%%%%%%%%%%%%%%%%%%%%%%%%%%%%%%%%%%%%%

A confidence interval typically describes uncertainty in an estimated
parameter or mean.

A prediction interval describes uncertainty in a future observation or
realization.

These are not generally the same width.

Prediction intervals usually include additional variability from the
system itself.

%%%%%%%%%%%%%%%%%%%%%%%%%%%%%%%%%%%%%%%%%%%%%%%%%%%%%%%%%%%%
\subsection{Rare-Event Simulation}
\label{subsec:simulation-rare-events}
%%%%%%%%%%%%%%%%%%%%%%%%%%%%%%%%%%%%%%%%%%%%%%%%%%%%%%%%%%%%

Some applications care primarily about

\[
\boxed{
P(Y>c),
}
\]

where the event is extremely unlikely.

Examples include:

\[
\boxed{
\text{structural failure},
}
\]

\[
\boxed{
\text{extreme financial loss},
}
\]

and

\[
\boxed{
\text{system overload}.
}
\]

Naive simulation may require an enormous number of replications.

%%%%%%%%%%%%%%%%%%%%%%%%%%%%%%%%%%%%%%%%%%%%%%%%%%%%%%%%%%%%
\subsection{Rare-Event Variance Reduction}
\label{subsec:simulation-rare-event-variance-reduction}
%%%%%%%%%%%%%%%%%%%%%%%%%%%%%%%%%%%%%%%%%%%%%%%%%%%%%%%%%%%%

Possible techniques include:

\[
\boxed{
\text{importance sampling},
}
\]

\[
\boxed{
\text{splitting methods},
}
\]

and specialized rare-event algorithms.

The goal is to observe the important region more frequently while
correcting for altered sampling.

%%%%%%%%%%%%%%%%%%%%%%%%%%%%%%%%%%%%%%%%%%%%%%%%%%%%%%%%%%%%
\subsection{Simulation-Based Optimization}
\label{subsec:simulation-based-optimization}
%%%%%%%%%%%%%%%%%%%%%%%%%%%%%%%%%%%%%%%%%%%%%%%%%%%%%%%%%%%%

Suppose decision variable \(x\) controls a simulation.

The objective may be

\[
\boxed{
J(x)
=
\mathbb{E}
[
C(x,\xi)
].
}
\]

If \(J(x)\) has no analytic form, it can be estimated by simulation:

\[
\boxed{
\widehat{J}_N(x)
=
\frac1N
\sum_{i=1}^{N}
C(x,\xi_i).
}
\]

An optimization algorithm then searches over \(x\).

%%%%%%%%%%%%%%%%%%%%%%%%%%%%%%%%%%%%%%%%%%%%%%%%%%%%%%%%%%%%
\subsection{Noisy Objective Functions}
\label{subsec:noisy-simulation-objectives}
%%%%%%%%%%%%%%%%%%%%%%%%%%%%%%%%%%%%%%%%%%%%%%%%%%%%%%%%%%%%

Simulation-based objectives may contain sampling noise.

Thus two evaluations at the same \(x\) can differ:

\[
\boxed{
\widehat{J}^{(1)}(x)
\neq
\widehat{J}^{(2)}(x).
}
\]

Optimization methods must account for this variability.

Possible approaches include:

\[
\boxed{
\text{larger sample sizes},
}
\]

\[
\boxed{
\text{common random numbers},
}
\]

or

\[
\boxed{
\text{stochastic optimization methods}.
}
\]

%%%%%%%%%%%%%%%%%%%%%%%%%%%%%%%%%%%%%%%%%%%%%%%%%%%%%%%%%%%%
\subsection{Common Random Numbers in Simulation Optimization}
\label{subsec:simulation-common-random-numbers}
%%%%%%%%%%%%%%%%%%%%%%%%%%%%%%%%%%%%%%%%%%%%%%%%%%%%%%%%%%%%

Suppose decisions \(x_A\) and \(x_B\) are compared.

Use identical random inputs

\[
\xi_1,\ldots,\xi_N
\]

for both.

Then compare

\[
\boxed{
\widehat{J}(x_A)
-
\widehat{J}(x_B).
}
\]

Correlation induced by common randomness can reduce variance in the
difference.

%%%%%%%%%%%%%%%%%%%%%%%%%%%%%%%%%%%%%%%%%%%%%%%%%%%%%%%%%%%%
\subsection{Surrogate Models}
\label{subsec:simulation-surrogate-models}
%%%%%%%%%%%%%%%%%%%%%%%%%%%%%%%%%%%%%%%%%%%%%%%%%%%%%%%%%%%%

If a simulation

\[
M(\theta)
\]

is expensive, approximate it using a cheaper model

\[
\boxed{
\widetilde{M}(\theta).
}
\]

The surrogate may be built from simulation samples.

Applications include:

\[
\boxed{
\text{optimization},
}
\]

\[
\boxed{
\text{uncertainty analysis},
}
\]

and

\[
\boxed{
\text{interactive exploration}.
}
\]

%%%%%%%%%%%%%%%%%%%%%%%%%%%%%%%%%%%%%%%%%%%%%%%%%%%%%%%%%%%%
\subsection{Response Surfaces}
\label{subsec:response-surfaces}
%%%%%%%%%%%%%%%%%%%%%%%%%%%%%%%%%%%%%%%%%%%%%%%%%%%%%%%%%%%%

A response surface approximates a quantity of interest as a function of
parameters:

\[
\boxed{
Q(\theta)
\approx
\widetilde{Q}(\theta).
}
\]

A simple response surface may use a polynomial:

\[
\boxed{
\widetilde{Q}(\theta)
=
\beta_0
+
\sum_i
\beta_i\theta_i
+
\sum_{i,j}
\beta_{ij}\theta_i\theta_j.
}
\]

Higher-order or nonparametric surrogates are also possible.

%%%%%%%%%%%%%%%%%%%%%%%%%%%%%%%%%%%%%%%%%%%%%%%%%%%%%%%%%%%%
\subsection{Surrogate Validation}
\label{subsec:surrogate-validation}
%%%%%%%%%%%%%%%%%%%%%%%%%%%%%%%%%%%%%%%%%%%%%%%%%%%%%%%%%%%%

A surrogate introduces approximation error.

Therefore compare surrogate predictions with independent full simulations.

A surrogate should not be trusted merely because it fits its training
points well.

%%%%%%%%%%%%%%%%%%%%%%%%%%%%%%%%%%%%%%%%%%%%%%%%%%%%%%%%%%%%
\subsection{Reduced-Order Simulation}
\label{subsec:reduced-order-simulation}
%%%%%%%%%%%%%%%%%%%%%%%%%%%%%%%%%%%%%%%%%%%%%%%%%%%%%%%%%%%%

Suppose the original state satisfies

\[
x(t)
\in
\mathbb{R}^N.
\]

Approximate

\[
\boxed{
x(t)
\approx
V_r a(t),
}
\]

where

\[
r\ll N.
\]

Then simulate the reduced coordinates

\[
a(t)
\]

instead of all \(N\) original variables.

This can greatly reduce computation.

%%%%%%%%%%%%%%%%%%%%%%%%%%%%%%%%%%%%%%%%%%%%%%%%%%%%%%%%%%%%
\subsection{Model Reduction Tradeoff}
\label{subsec:model-reduction-tradeoff}
%%%%%%%%%%%%%%%%%%%%%%%%%%%%%%%%%%%%%%%%%%%%%%%%%%%%%%%%%%%%

Reduced-order models exchange accuracy for speed.

The central question is whether the reduced model preserves the quantities
of interest.

Validation must be performed over the parameter and state ranges where the
reduced model will be used.

%%%%%%%%%%%%%%%%%%%%%%%%%%%%%%%%%%%%%%%%%%%%%%%%%%%%%%%%%%%%
\subsection{Digital Experiments}
\label{subsec:digital-experiments}
%%%%%%%%%%%%%%%%%%%%%%%%%%%%%%%%%%%%%%%%%%%%%%%%%%%%%%%%%%%%

A simulation study should be designed like an experiment.

Specify:

\[
\boxed{
\text{factors},
}
\]

\[
\boxed{
\text{factor levels},
}
\]

\[
\boxed{
\text{responses},
}
\]

and

\[
\boxed{
\text{replications}.
}
\]

Statistical design-of-experiments principles can improve the efficiency of
simulation studies.

%%%%%%%%%%%%%%%%%%%%%%%%%%%%%%%%%%%%%%%%%%%%%%%%%%%%%%%%%%%%
\subsection{Factorial Simulation Designs}
\label{subsec:factorial-simulation-design}
%%%%%%%%%%%%%%%%%%%%%%%%%%%%%%%%%%%%%%%%%%%%%%%%%%%%%%%%%%%%

Suppose two parameters each have two selected levels.

A full factorial design contains

\[
\boxed{
2^2=4
}
\]

parameter combinations.

For \(k\) two-level factors, the design contains

\[
\boxed{
2^k
}
\]

combinations.

Factorial designs reveal interactions that one-at-a-time sweeps may miss.

%%%%%%%%%%%%%%%%%%%%%%%%%%%%%%%%%%%%%%%%%%%%%%%%%%%%%%%%%%%%
\subsection{Replication and Randomization}
\label{subsec:simulation-replication-randomization}
%%%%%%%%%%%%%%%%%%%%%%%%%%%%%%%%%%%%%%%%%%%%%%%%%%%%%%%%%%%%

Replication estimates stochastic variability.

Randomization prevents systematic ordering effects from contaminating
comparisons.

Although computer simulations are controlled experiments, random execution
order can still help when computational environments or shared resources
vary over time.

%%%%%%%%%%%%%%%%%%%%%%%%%%%%%%%%%%%%%%%%%%%%%%%%%%%%%%%%%%%%
\subsection{Simulation Data Management}
\label{subsec:simulation-data-management}
%%%%%%%%%%%%%%%%%%%%%%%%%%%%%%%%%%%%%%%%%%%%%%%%%%%%%%%%%%%%

Large simulation studies may produce:

\[
\boxed{
\text{input configurations},
}
\]

\[
\boxed{
\text{raw trajectories},
}
\]

\[
\boxed{
\text{summary statistics},
}
\]

\[
\boxed{
\text{checkpoints},
}
\]

and

\[
\boxed{
\text{plots}.
}
\]

Data organization should preserve the relationship between each result and
the simulation settings that generated it.

%%%%%%%%%%%%%%%%%%%%%%%%%%%%%%%%%%%%%%%%%%%%%%%%%%%%%%%%%%%%
\subsection{Simulation Metadata}
\label{subsec:simulation-metadata}
%%%%%%%%%%%%%%%%%%%%%%%%%%%%%%%%%%%%%%%%%%%%%%%%%%%%%%%%%%%%

Each simulation result should ideally record:

\[
\boxed{
\text{parameter values},
}
\]

\[
\boxed{
\text{initial conditions},
}
\]

\[
\boxed{
\text{numerical method},
}
\]

\[
\boxed{
\text{step size or tolerances},
}
\]

\[
\boxed{
\text{random seed},
}
\]

and

\[
\boxed{
\text{software version}.
}
\]

Without metadata, large simulation studies become difficult to interpret
or reproduce.

%%%%%%%%%%%%%%%%%%%%%%%%%%%%%%%%%%%%%%%%%%%%%%%%%%%%%%%%%%%%
\subsection{Checkpointing and Restart}
\label{subsec:simulation-checkpointing}
%%%%%%%%%%%%%%%%%%%%%%%%%%%%%%%%%%%%%%%%%%%%%%%%%%%%%%%%%%%%

Long simulations should periodically save the current state.

A checkpoint may contain:

\[
\boxed{
x(t_k),
}
\]

\[
\boxed{
t_k,
}
\]

\[
\boxed{
\text{random generator state},
}
\]

and

\[
\boxed{
\text{solver state}.
}
\]

The simulation can then resume without restarting from \(t=0\).

%%%%%%%%%%%%%%%%%%%%%%%%%%%%%%%%%%%%%%%%%%%%%%%%%%%%%%%%%%%%
\subsection{Parallel Simulation}
\label{subsec:parallel-simulation}
%%%%%%%%%%%%%%%%%%%%%%%%%%%%%%%%%%%%%%%%%%%%%%%%%%%%%%%%%%%%

Many simulation tasks are embarrassingly parallel.

Examples include:

\[
\boxed{
\text{independent stochastic replications},
}
\]

\[
\boxed{
\text{parameter combinations},
}
\]

and

\[
\boxed{
\text{scenario runs}.
}
\]

If each simulation is independent, the work can be distributed across many
processors with little communication.

%%%%%%%%%%%%%%%%%%%%%%%%%%%%%%%%%%%%%%%%%%%%%%%%%%%%%%%%%%%%
\subsection{Parallel Ensemble Average}
\label{subsec:parallel-ensemble-average}
%%%%%%%%%%%%%%%%%%%%%%%%%%%%%%%%%%%%%%%%%%%%%%%%%%%%%%%%%%%%

Suppose processor \(j\) computes local sum

\[
\boxed{
S_j
=
\sum_{r\in I_j}
Y_r.
}
\]

After all workers finish, perform a reduction:

\[
\boxed{
\overline{Y}
=
\frac{
\sum_jS_j
}{
N
}.
}
\]

Only the final summary requires communication.

%%%%%%%%%%%%%%%%%%%%%%%%%%%%%%%%%%%%%%%%%%%%%%%%%%%%%%%%%%%%
\subsection{Random Streams in Parallel Simulation}
\label{subsec:parallel-simulation-random-streams}
%%%%%%%%%%%%%%%%%%%%%%%%%%%%%%%%%%%%%%%%%%%%%%%%%%%%%%%%%%%%

Parallel stochastic simulations should use carefully managed random-number
streams.

Different workers should not accidentally generate overlapping sequences.

Reproducible stream assignment also makes parallel experiments easier to
verify.

%%%%%%%%%%%%%%%%%%%%%%%%%%%%%%%%%%%%%%%%%%%%%%%%%%%%%%%%%%%%
\subsection{Visualization of Simulation Results}
\label{subsec:simulation-visualization}
%%%%%%%%%%%%%%%%%%%%%%%%%%%%%%%%%%%%%%%%%%%%%%%%%%%%%%%%%%%%

Useful visualizations include:

\[
\boxed{
\text{trajectories},
}
\]

\[
\boxed{
\text{phase portraits},
}
\]

\[
\boxed{
\text{heat maps},
}
\]

\[
\boxed{
\text{histograms},
}
\]

\[
\boxed{
\text{quantile bands},
}
\]

and

\[
\boxed{
\text{animations}.
}
\]

The visualization should correspond to the scientific quantity of
interest.

%%%%%%%%%%%%%%%%%%%%%%%%%%%%%%%%%%%%%%%%%%%%%%%%%%%%%%%%%%%%
\subsection{Phase Portraits}
\label{subsec:simulation-phase-portraits}
%%%%%%%%%%%%%%%%%%%%%%%%%%%%%%%%%%%%%%%%%%%%%%%%%%%%%%%%%%%%

For a two-dimensional state

\[
(x(t),y(t)),
\]

plot

\[
\boxed{
y
\text{ versus }
x
}
\]

rather than plotting both only against time.

Phase portraits reveal:

\[
\boxed{
\text{equilibria},
}
\]

\[
\boxed{
\text{cycles},
}
\]

and

\[
\boxed{
\text{qualitative trajectory structure}.
}
\]

%%%%%%%%%%%%%%%%%%%%%%%%%%%%%%%%%%%%%%%%%%%%%%%%%%%%%%%%%%%%
\subsection{Ensemble Visualization}
\label{subsec:ensemble-visualization}
%%%%%%%%%%%%%%%%%%%%%%%%%%%%%%%%%%%%%%%%%%%%%%%%%%%%%%%%%%%%

Plotting hundreds of stochastic trajectories simultaneously can become
unreadable.

Instead, summarize with:

\[
\boxed{
\text{median trajectory},
}
\]

and quantile envelopes such as

\[
\boxed{
5\%,
25\%,
75\%,
95\%.
}
\]

This displays uncertainty without overwhelming the figure.

%%%%%%%%%%%%%%%%%%%%%%%%%%%%%%%%%%%%%%%%%%%%%%%%%%%%%%%%%%%%
\subsection{Visual Plausibility Is Not Validation}
\label{subsec:simulation-visual-plausibility}
%%%%%%%%%%%%%%%%%%%%%%%%%%%%%%%%%%%%%%%%%%%%%%%%%%%%%%%%%%%%

An animation may appear realistic while the model is mathematically or
physically wrong.

Simulation quality must be assessed using:

\[
\boxed{
\text{convergence},
}
\]

\[
\boxed{
\text{conservation},
}
\]

\[
\boxed{
\text{data comparison},
}
\]

and

\[
\boxed{
\text{uncertainty analysis}.
}
\]

A visually appealing result is not sufficient evidence.

%%%%%%%%%%%%%%%%%%%%%%%%%%%%%%%%%%%%%%%%%%%%%%%%%%%%%%%%%%%%
\subsection{Example: Projectile Simulation}
\label{subsec:projectile-simulation}
%%%%%%%%%%%%%%%%%%%%%%%%%%%%%%%%%%%%%%%%%%%%%%%%%%%%%%%%%%%%

Ignoring air resistance, projectile motion satisfies

\[
\boxed{
x''=0,
}
\]

\[
\boxed{
y''=-g.
}
\]

Introduce velocities

\[
v_x=x',
\qquad
v_y=y'.
\]

Then

\[
\boxed{
x'=v_x,
}
\]

\[
\boxed{
y'=v_y,
}
\]

\[
\boxed{
v_x'=0,
}
\]

and

\[
\boxed{
v_y'=-g.
}
\]

A numerical ODE solver can simulate the trajectory until

\[
y(t)=0.
\]

%%%%%%%%%%%%%%%%%%%%%%%%%%%%%%%%%%%%%%%%%%%%%%%%%%%%%%%%%%%%
\subsection{Projectile Event Detection}
\label{subsec:projectile-event-detection}
%%%%%%%%%%%%%%%%%%%%%%%%%%%%%%%%%%%%%%%%%%%%%%%%%%%%%%%%%%%%

Define event function

\[
\boxed{
g_{\mathrm{event}}(t,x,y,v_x,v_y)
=
y.
}
\]

The impact occurs when

\[
\boxed{
y=0
}
\]

while the projectile is moving downward.

Accurate event location avoids relying on whichever discrete step first
produces a negative height.

%%%%%%%%%%%%%%%%%%%%%%%%%%%%%%%%%%%%%%%%%%%%%%%%%%%%%%%%%%%%
\subsection{Example: Predator--Prey Simulation}
\label{subsec:predator-prey-simulation}
%%%%%%%%%%%%%%%%%%%%%%%%%%%%%%%%%%%%%%%%%%%%%%%%%%%%%%%%%%%%

The Lotka--Volterra equations are

\[
\boxed{
\frac{dx}{dt}
=
\alpha x-\beta xy,
}
\]

\[
\boxed{
\frac{dy}{dt}
=
\delta xy-\gamma y.
}
\]

Here:

\[
x(t)
=
\text{prey population},
\]

and

\[
y(t)
=
\text{predator population}.
\]

Simulation can reveal oscillatory population dynamics and dependence on
parameters.

%%%%%%%%%%%%%%%%%%%%%%%%%%%%%%%%%%%%%%%%%%%%%%%%%%%%%%%%%%%%
\subsection{Example: Epidemic Simulation}
\label{subsec:epidemic-simulation}
%%%%%%%%%%%%%%%%%%%%%%%%%%%%%%%%%%%%%%%%%%%%%%%%%%%%%%%%%%%%

An SEIR-type model may use

\[
\boxed{
S+E+I+R=N.
}
\]

A possible deterministic system is

\[
\boxed{
\frac{dS}{dt}
=
-\beta
\frac{SI}{N},
}
\]

\[
\boxed{
\frac{dE}{dt}
=
\beta
\frac{SI}{N}
-
\sigma E,
}
\]

\[
\boxed{
\frac{dI}{dt}
=
\sigma E
-
\gamma I,
}
\]

\[
\boxed{
\frac{dR}{dt}
=
\gamma I.
}
\]

Simulation can compare interventions by modifying parameters or movement
rules.

%%%%%%%%%%%%%%%%%%%%%%%%%%%%%%%%%%%%%%%%%%%%%%%%%%%%%%%%%%%%
\subsection{Example: Inventory Simulation}
\label{subsec:inventory-simulation}
%%%%%%%%%%%%%%%%%%%%%%%%%%%%%%%%%%%%%%%%%%%%%%%%%%%%%%%%%%%%

Let

\[
I_t
\]

denote inventory.

A simple discrete update is

\[
\boxed{
I_{t+1}
=
I_t
+
Q_t
-
D_t,
}
\]

where

\[
Q_t
\]

is replenishment and

\[
D_t
\]

is demand.

If demand is random, repeated simulation estimates:

\[
\boxed{
\text{stockout probability},
}
\]

\[
\boxed{
\text{holding cost},
}
\]

and

\[
\boxed{
\text{service level}.
}
\]

%%%%%%%%%%%%%%%%%%%%%%%%%%%%%%%%%%%%%%%%%%%%%%%%%%%%%%%%%%%%
\subsection{Example: Reliability Simulation}
\label{subsec:reliability-simulation}
%%%%%%%%%%%%%%%%%%%%%%%%%%%%%%%%%%%%%%%%%%%%%%%%%%%%%%%%%%%%

Suppose component lifetimes are random:

\[
\boxed{
T_i.
}
\]

System failure time

\[
T_{\mathrm{sys}}
\]

depends on how components are connected.

Monte Carlo simulation samples the component lifetimes, evaluates the
system failure rule, and estimates quantities such as

\[
\boxed{
P(T_{\mathrm{sys}}\leq t).
}
\]

%%%%%%%%%%%%%%%%%%%%%%%%%%%%%%%%%%%%%%%%%%%%%%%%%%%%%%%%%%%%
\subsection{Example: Financial Path Simulation}
\label{subsec:financial-path-simulation}
%%%%%%%%%%%%%%%%%%%%%%%%%%%%%%%%%%%%%%%%%%%%%%%%%%%%%%%%%%%%

A stochastic asset model may satisfy

\[
\boxed{
dS_t
=
\mu S_t\,dt
+
\sigma S_t\,dW_t.
}
\]

Simulated paths

\[
S_t^{(1)},
\ldots,
S_t^{(N)}
\]

can estimate distributions of future asset values or path-dependent
payoffs.

The simulation must account for both stochastic sampling error and
time-discretization error.

%%%%%%%%%%%%%%%%%%%%%%%%%%%%%%%%%%%%%%%%%%%%%%%%%%%%%%%%%%%%
\subsection{Simulation Credibility}
\label{subsec:simulation-credibility}
%%%%%%%%%%%%%%%%%%%%%%%%%%%%%%%%%%%%%%%%%%%%%%%%%%%%%%%%%%%%

A credible simulation should provide evidence concerning:

\[
\boxed{
\text{model assumptions},
}
\]

\[
\boxed{
\text{implementation correctness},
}
\]

\[
\boxed{
\text{numerical accuracy},
}
\]

\[
\boxed{
\text{parameter uncertainty},
}
\]

and

\[
\boxed{
\text{validation against observations}.
}
\]

Credibility is built through multiple independent checks.

%%%%%%%%%%%%%%%%%%%%%%%%%%%%%%%%%%%%%%%%%%%%%%%%%%%%%%%%%%%%
\subsection{Simulation Workflow}
\label{subsec:mathematical-simulation-workflow}
%%%%%%%%%%%%%%%%%%%%%%%%%%%%%%%%%%%%%%%%%%%%%%%%%%%%%%%%%%%%

A complete simulation workflow is:

\begin{enumerate}
    \item define the scientific or operational question;
    \item identify quantities of interest;
    \item formulate the mathematical model;
    \item identify states, parameters, inputs, and outputs;
    \item state assumptions explicitly;
    \item determine whether the model is deterministic or stochastic;
    \item determine whether it is continuous-time, discrete-time, or
          event-driven;
    \item specify initial and boundary conditions;
    \item select parameter values or parameter distributions;
    \item choose numerical methods;
    \item select time step, mesh, or solver tolerances;
    \item implement the simulation;
    \item verify the implementation on simple cases;
    \item monitor invariants and constraints;
    \item perform numerical convergence studies;
    \item determine the appropriate simulation horizon;
    \item identify and handle transient behavior;
    \item perform stochastic replications when required;
    \item estimate simulation uncertainty;
    \item perform sensitivity analysis;
    \item calibrate uncertain parameters when data are available;
    \item validate against independent observations;
    \item compare alternative scenarios;
    \item document metadata and random seeds;
    \item visualize important outputs;
    \item report uncertainty and limitations;
    \item preserve enough information for reproducibility.
\end{enumerate}

%%%%%%%%%%%%%%%%%%%%%%%%%%%%%%%%%%%%%%%%%%%%%%%%%%%%%%%%%%%%
\subsection{Choosing a Simulation Approach}
\label{subsec:choosing-simulation-approach}
%%%%%%%%%%%%%%%%%%%%%%%%%%%%%%%%%%%%%%%%%%%%%%%%%%%%%%%%%%%%

For smooth deterministic dynamics, consider:

\[
\boxed{
\text{ODE or PDE simulation}.
}
\]

For systems driven by discrete events, consider:

\[
\boxed{
\text{discrete-event simulation}.
}
\]

For heterogeneous interacting individuals, consider:

\[
\boxed{
\text{agent-based simulation}.
}
\]

For random systems with known transition structure, consider:

\[
\boxed{
\text{stochastic-process simulation}.
}
\]

For uncertainty propagation, consider:

\[
\boxed{
\text{Monte Carlo ensembles}.
}
\]

%%%%%%%%%%%%%%%%%%%%%%%%%%%%%%%%%%%%%%%%%%%%%%%%%%%%%%%%%%%%
\subsection{Common Errors}
\label{subsec:mathematical-simulation-common-errors}
%%%%%%%%%%%%%%%%%%%%%%%%%%%%%%%%%%%%%%%%%%%%%%%%%%%%%%%%%%%%

\subsubsection{Confusing the Model With Reality}

A simulation approximates a mathematical model.

Even a perfectly computed solution may fail to represent the real system
adequately.

%%%%%%%%%%%%%%%%%%%%%%%%%%%%%%%%%%%%%%%%%%%%%%%%%%%%%%%%%%%%
\subsubsection{Confusing the Model With the Numerical Method}

For example,

\[
\dot{x}=f(x)
\]

is the model.

Euler's method is only one possible numerical algorithm for simulating it.

Changing the numerical solver should not change the mathematical model.

%%%%%%%%%%%%%%%%%%%%%%%%%%%%%%%%%%%%%%%%%%%%%%%%%%%%%%%%%%%%
\subsubsection{Using One Stochastic Run as the Expected Outcome}

One realization may be highly atypical.

Use independent replications and statistical summaries.

%%%%%%%%%%%%%%%%%%%%%%%%%%%%%%%%%%%%%%%%%%%%%%%%%%%%%%%%%%%%
\subsubsection{Reporting Only the Mean}

A mean can hide:

\[
\boxed{
\text{large variance},
}
\]

\[
\boxed{
\text{skewness},
}
\]

\[
\boxed{
\text{multimodality},
}
\]

and

\[
\boxed{
\text{extreme tails}.
}
\]

Report distributional information when it matters.

%%%%%%%%%%%%%%%%%%%%%%%%%%%%%%%%%%%%%%%%%%%%%%%%%%%%%%%%%%%%
\subsubsection{Ignoring Warm-Up Bias}

A queue initialized empty may not represent long-run behavior.

Steady-state estimation must account for transient initialization effects.

%%%%%%%%%%%%%%%%%%%%%%%%%%%%%%%%%%%%%%%%%%%%%%%%%%%%%%%%%%%%
\subsubsection{Discarding Warm-Up in a Finite-Horizon Problem}

If the actual system begins empty, the transient may be physically
meaningful.

Do not remove it merely because steady-state simulation often uses
warm-up deletion.

%%%%%%%%%%%%%%%%%%%%%%%%%%%%%%%%%%%%%%%%%%%%%%%%%%%%%%%%%%%%
\subsubsection{Using an Arbitrary Simulation Horizon}

A horizon that is too short can miss important long-term behavior.

A horizon that is unnecessarily long wastes computation.

Choose it according to the scientific objective and system time scales.

%%%%%%%%%%%%%%%%%%%%%%%%%%%%%%%%%%%%%%%%%%%%%%%%%%%%%%%%%%%%
\subsubsection{Ignoring Numerical Stability}

An exploding simulation may result from an unstable discretization rather
than an unstable physical model.

Repeat with a smaller time step or a more appropriate method.

%%%%%%%%%%%%%%%%%%%%%%%%%%%%%%%%%%%%%%%%%%%%%%%%%%%%%%%%%%%%
\subsubsection{Ignoring Conservation Laws}

If the exact model conserves mass or population, major numerical drift is
a warning sign.

Track known invariants.

%%%%%%%%%%%%%%%%%%%%%%%%%%%%%%%%%%%%%%%%%%%%%%%%%%%%%%%%%%%%
\subsubsection{Treating Tiny Negative Values as Automatically Physical}

A concentration of

\[
-10^{-8}
\]

may simply be numerical error.

The model's admissible state space should guide interpretation.

%%%%%%%%%%%%%%%%%%%%%%%%%%%%%%%%%%%%%%%%%%%%%%%%%%%%%%%%%%%%
\subsubsection{Changing Many Parameters at Once Without Experimental Design}

The resulting output change may be difficult to attribute.

Use structured scenario analysis or experimental design.

%%%%%%%%%%%%%%%%%%%%%%%%%%%%%%%%%%%%%%%%%%%%%%%%%%%%%%%%%%%%
\subsubsection{Using One-at-a-Time Sensitivity to Infer Interactions}

One-at-a-time methods do not reveal parameter interactions.

Use global or factorial approaches when interactions matter.

%%%%%%%%%%%%%%%%%%%%%%%%%%%%%%%%%%%%%%%%%%%%%%%%%%%%%%%%%%%%
\subsubsection{Calibrating and Validating on the Same Data}

A model can appear artificially accurate when tested on the data used to
fit it.

Independent validation data provide stronger evidence.

%%%%%%%%%%%%%%%%%%%%%%%%%%%%%%%%%%%%%%%%%%%%%%%%%%%%%%%%%%%%
\subsubsection{Assuming a Good Fit Proves the Model Mechanism Is Correct}

Different parameter sets or even different models may fit the same data.

Predictive agreement does not automatically identify the true mechanism.

%%%%%%%%%%%%%%%%%%%%%%%%%%%%%%%%%%%%%%%%%%%%%%%%%%%%%%%%%%%%
\subsubsection{Ignoring Identifiability}

An optimizer may return one parameter vector even when many nearly
equivalent vectors fit the data.

Investigate identifiability and parameter uncertainty.

%%%%%%%%%%%%%%%%%%%%%%%%%%%%%%%%%%%%%%%%%%%%%%%%%%%%%%%%%%%%
\subsubsection{Assuming More Simulation Runs Fix Model Bias}

Increasing replication reduces Monte Carlo error.

It does not eliminate:

\[
\boxed{
\text{model-form error}
}
\]

or

\[
\boxed{
\text{systematic parameter bias}.
}
\]

%%%%%%%%%%%%%%%%%%%%%%%%%%%%%%%%%%%%%%%%%%%%%%%%%%%%%%%%%%%%
\subsubsection{Using a Surrogate Outside Its Training Region}

A surrogate may interpolate well but extrapolate poorly.

Validate over the region where it will actually be used.

%%%%%%%%%%%%%%%%%%%%%%%%%%%%%%%%%%%%%%%%%%%%%%%%%%%%%%%%%%%%
\subsubsection{Saving Results Without Configuration Metadata}

A trajectory without parameter values, solver tolerances, and seed
information may be impossible to reproduce.

Store simulation metadata with outputs.

%%%%%%%%%%%%%%%%%%%%%%%%%%%%%%%%%%%%%%%%%%%%%%%%%%%%%%%%%%%%
\subsubsection{Assuming an Animation Is Evidence of Correctness}

A realistic-looking animation can still arise from incorrect equations or
unstable numerics.

Use quantitative verification and validation.

%%%%%%%%%%%%%%%%%%%%%%%%%%%%%%%%%%%%%%%%%%%%%%%%%%%%%%%%%%%%
\subsection{Applications}
\label{subsec:mathematical-simulation-applications}
%%%%%%%%%%%%%%%%%%%%%%%%%%%%%%%%%%%%%%%%%%%%%%%%%%%%%%%%%%%%

\begin{applicationbox}

\textbf{Epidemiology.}
Compartment and agent-based models simulate infection spread, intervention
strategies, mobility, vaccination, and healthcare demand.

\medskip

\textbf{Ecology.}
Population, predator--prey, competition, dispersal, and spatial models can
be explored computationally.

\medskip

\textbf{Engineering.}
Simulation predicts stress, temperature, fluid flow, control-system
response, reliability, and manufacturing performance.

\medskip

\textbf{Physics.}
Particle motion, waves, fields, statistical systems, and dynamical systems
are commonly studied through numerical simulation.

\medskip

\textbf{Operations Research.}
Queues, inventories, transportation systems, staffing policies, and
service operations are evaluated through discrete-event and stochastic
simulation.

\medskip

\textbf{Finance.}
Stochastic path simulation is used for pricing, risk analysis, portfolio
evaluation, and scenario generation.

\medskip

\textbf{Environmental Science.}
Climate, hydrology, pollution transport, ecosystems, and ocean-atmosphere
systems require large-scale mathematical simulations.

\medskip

\textbf{Mathematical Research.}
Simulation provides numerical experiments that can reveal patterns,
suggest conjectures, test asymptotic regimes, and explore models before
formal analysis is available.

\end{applicationbox}

%%%%%%%%%%%%%%%%%%%%%%%%%%%%%%%%%%%%%%%%%%%%%%%%%%%%%%%%%%%%
\subsection{Exercises}
\label{subsec:mathematical-simulation-exercises}
%%%%%%%%%%%%%%%%%%%%%%%%%%%%%%%%%%%%%%%%%%%%%%%%%%%%%%%%%%%%

\begin{exercise}[1]
Define mathematical simulation.
\end{exercise}

\begin{exercise}[1]
Distinguish a model from a simulation.
\end{exercise}

\begin{exercise}[1]
Define a state variable.
\end{exercise}

\begin{exercise}[1]
Define a simulation parameter.
\end{exercise}

\begin{exercise}[1]
Distinguish deterministic and stochastic simulation.
\end{exercise}

\begin{exercise}[1]
Distinguish static and dynamic simulation.
\end{exercise}

\begin{exercise}[1]
Define event-driven simulation.
\end{exercise}

\begin{exercise}[1]
Define agent-based simulation.
\end{exercise}

\begin{exercise}[1]
Define a simulation replication.
\end{exercise}

\begin{exercise}[1]
Define model validation.
\end{exercise}

\begin{exercise}[2]
Identify the states and parameters in a logistic population model.
\end{exercise}

\begin{exercise}[2]
Perform three updates of a discrete logistic model.
\end{exercise}

\begin{exercise}[2]
Write the update rule for a simple inventory model.
\end{exercise}

\begin{exercise}[2]
Describe the state variables of a single-server queue.
\end{exercise}

\begin{exercise}[3]
Construct a future-event list containing three arrivals and two
departures.
\end{exercise}

\begin{exercise}[3]
Process one event and update the event calendar.
\end{exercise}

\begin{exercise}[3]
Compare time-driven and event-driven simulation for a queue.
\end{exercise}

\begin{exercise}[3]
Construct a simple synchronous agent update rule.
\end{exercise}

\begin{exercise}[3]
Explain how asynchronous updating could change the outcome.
\end{exercise}

\begin{exercise}[2]
Write the SIR model and verify that

\[
S+I+R
\]

is conserved.
\end{exercise}

\begin{exercise}[3]
Design a numerical diagnostic for population conservation.
\end{exercise}

\begin{exercise}[3]
Construct a network-state model using an adjacency matrix.
\end{exercise}

\begin{exercise}[3]
Write a two-state Markov-chain simulation algorithm.
\end{exercise}

\begin{exercise}[2]
Choose a reasonable simulation horizon for a specified finite-duration
system and explain the choice.
\end{exercise}

\begin{exercise}[3]
Explain why an empty initial queue may bias steady-state estimates.
\end{exercise}

\begin{exercise}[3]
Explain when a warm-up period should not be discarded.
\end{exercise}

\begin{exercise}[3]
Design a time-step refinement study for a dynamic simulation.
\end{exercise}

\begin{exercise}[3]
Explain how to distinguish physical instability from numerical
instability.
\end{exercise}

\begin{exercise}[3]
Construct a positivity check for a compartment model.
\end{exercise}

\begin{exercise}[2]
Compute the ensemble mean from five stochastic replications.
\end{exercise}

\begin{exercise}[2]
Compute the sample variance across replications.
\end{exercise}

\begin{exercise}[3]
Construct an approximate confidence interval for a simulation mean.
\end{exercise}

\begin{exercise}[3]
Explain why a pointwise confidence interval is not automatically a
simultaneous trajectory band.
\end{exercise}

\begin{exercise}[3]
Distinguish aleatoric and parametric uncertainty.
\end{exercise}

\begin{exercise}[3]
Design a Monte Carlo uncertainty-propagation experiment.
\end{exercise}

\begin{exercise}[3]
Estimate an exceedance probability

\[
P(Y>c)
\]

from simulation output.
\end{exercise}

\begin{exercise}[3]
Estimate a simulation quantile from sorted outputs.
\end{exercise}

\begin{exercise}[2]
Construct three explicit scenarios for a population model.
\end{exercise}

\begin{exercise}[3]
Design a controlled comparison between two policy scenarios.
\end{exercise}

\begin{exercise}[3]
Construct a two-parameter simulation sweep.
\end{exercise}

\begin{exercise}[3]
Approximate a local parameter sensitivity using centered differences.
\end{exercise}

\begin{exercise}[3]
Compute a normalized sensitivity coefficient.
\end{exercise}

\begin{exercise}[3]
Explain why one-at-a-time sensitivity misses interactions.
\end{exercise}

\begin{exercise}[3]
Describe global sensitivity analysis.
\end{exercise}

\begin{exercise}[2]
Construct a least-squares calibration objective.
\end{exercise}

\begin{exercise}[3]
Construct a weighted least-squares calibration objective.
\end{exercise}

\begin{exercise}[3]
Explain simulation-based parameter estimation.
\end{exercise}

\begin{exercise}[3]
Explain the relationship between calibration and inverse problems.
\end{exercise}

\begin{exercise}[3]
Give an example of a structurally unidentifiable model parameterization.
\end{exercise}

\begin{exercise}[3]
Describe signs of practical nonidentifiability.
\end{exercise}

\begin{exercise}[3]
Distinguish calibration, verification, and validation.
\end{exercise}

\begin{exercise}[3]
Design a benchmark test for a numerical simulation.
\end{exercise}

\begin{exercise}[3]
Construct a synthetic-data parameter-recovery experiment.
\end{exercise}

\begin{exercise}[3]
Perform a convergence-order calculation from supplied simulation errors.
\end{exercise}

\begin{exercise}[3]
Identify possible sources of discrepancy between simulation and
experiment.
\end{exercise}

\begin{exercise}[3]
Explain model-form error.
\end{exercise}

\begin{exercise}[3]
Explain why numerical refinement cannot eliminate model-form error.
\end{exercise}

\begin{exercise}[3]
Design a residual analysis for calibrated time-series data.
\end{exercise}

\begin{exercise}[3]
Explain the difference between confidence and prediction intervals.
\end{exercise}

\begin{exercise}[3]
Explain why naive simulation can be inefficient for rare-event
probabilities.
\end{exercise}

\begin{exercise}[3]
Describe an importance-sampling approach to rare-event simulation.
\end{exercise}

\begin{exercise}[3]
Construct a simulation-based optimization objective.
\end{exercise}

\begin{exercise}[3]
Explain why the objective may be noisy.
\end{exercise}

\begin{exercise}[3]
Explain how common random numbers can improve comparison of two decisions.
\end{exercise}

\begin{exercise}[3]
Construct a quadratic response-surface model for two parameters.
\end{exercise}

\begin{exercise}[3]
Describe how to validate a surrogate model.
\end{exercise}

\begin{exercise}[3]
Explain reduced-order simulation.
\end{exercise}

\begin{exercise}[3]
Describe the tradeoff between reduced-model speed and accuracy.
\end{exercise}

\begin{exercise}[3]
Construct a \(2^3\) factorial simulation design.
\end{exercise}

\begin{exercise}[3]
Explain why factorial designs reveal interactions.
\end{exercise}

\begin{exercise}[3]
List the metadata that should accompany a simulation result.
\end{exercise}

\begin{exercise}[3]
Design a checkpoint strategy for a long stochastic simulation.
\end{exercise}

\begin{exercise}[3]
Explain why random generator state may need to be checkpointed.
\end{exercise}

\begin{exercise}[3]
Design a parallel ensemble simulation.
\end{exercise}

\begin{exercise}[3]
Explain how to prevent overlapping random streams across workers.
\end{exercise}

\begin{exercise}[3]
Design an uncertainty visualization for a stochastic trajectory.
\end{exercise}

\begin{exercise}[3]
Explain why visualization alone is insufficient for validation.
\end{exercise}

\begin{exercise}[4]
Study discrete-event simulation theory for queueing systems.
\end{exercise}

\begin{exercise}[4]
Study stochastic simulation algorithms for continuous-time Markov chains.
\end{exercise}

\begin{exercise}[4]
Study exact stochastic chemical-reaction simulation.
\end{exercise}

\begin{exercise}[4]
Study agent-based epidemic models.
\end{exercise}

\begin{exercise}[4]
Study network epidemic simulation.
\end{exercise}

\begin{exercise}[4]
Study chaotic sensitivity and Lyapunov exponents in simulation.
\end{exercise}

\begin{exercise}[4]
Study structural identifiability of nonlinear ODE models.
\end{exercise}

\begin{exercise}[4]
Study profile likelihood for practical identifiability.
\end{exercise}

\begin{exercise}[4]
Study Bayesian calibration of simulation models.
\end{exercise}

\begin{exercise}[4]
Study model discrepancy in calibration.
\end{exercise}

\begin{exercise}[4]
Study global variance-based sensitivity indices.
\end{exercise}

\begin{exercise}[4]
Study uncertainty propagation using polynomial chaos.
\end{exercise}

\begin{exercise}[4]
Study rare-event splitting algorithms.
\end{exercise}

\begin{exercise}[4]
Study simulation optimization with noisy objectives.
\end{exercise}

\begin{exercise}[4]
Study surrogate-based optimization.
\end{exercise}

\begin{exercise}[4]
Study Gaussian-process surrogate modeling.
\end{exercise}

\begin{exercise}[4]
Study reduced-order modeling for dynamical systems.
\end{exercise}

\begin{exercise}[4]
Study verification and validation frameworks for computational models.
\end{exercise}

\begin{exercise}[4]
Design a complete reproducible simulation study from model formulation to
validation.
\end{exercise}

%%%%%%%%%%%%%%%%%%%%%%%%%%%%%%%%%%%%%%%%%%%%%%%%%%%%%%%%%%%%
\subsection{Conceptual Summary}
\label{subsec:mathematical-simulation-summary}
%%%%%%%%%%%%%%%%%%%%%%%%%%%%%%%%%%%%%%%%%%%%%%%%%%%%%%%%%%%%

Mathematical simulation combines

\[
\boxed{
\text{model}
+
\text{numerical method}
+
\text{computational experiment}.
}
\]

Simulations may be

\[
\boxed{
\text{deterministic or stochastic},
}
\]

\[
\boxed{
\text{static or dynamic},
}
\]

and

\[
\boxed{
\text{time-driven or event-driven}.
}
\]

A simulation state may evolve through

\[
\boxed{
\dot{x}
=
f(x,t;\theta)
}
\]

or

\[
\boxed{
x_{n+1}
=
F(x_n,n;\theta).
}
\]

Reliable simulation requires

\[
\boxed{
\text{verification}
+
\text{validation}
+
\text{uncertainty quantification}
+
\text{sensitivity analysis}.
}
\]

Calibration may be written as

\[
\boxed{
\min_\theta
J(\theta),
}
\]

where evaluating \(J\) may itself require a full simulation.

Stochastic studies require repeated runs and statistical summaries.

Simulation is therefore not merely running code. It is a controlled
mathematical experiment requiring careful modeling, numerical analysis,
statistical reasoning, and computational documentation.

%%%%%%%%%%%%%%%%%%%%%%%%%%%%%%%%%%%%%%%%%%%%%%%%%%%%%%%%%%%%
\subsection{Section Connections}
\label{subsec:mathematical-simulation-connections}
%%%%%%%%%%%%%%%%%%%%%%%%%%%%%%%%%%%%%%%%%%%%%%%%%%%%%%%%%%%%

\chapterconnection{
Mathematical Simulation combines many of the computational tools developed
throughout Chapter~11. Numerical ODE and PDE methods generate deterministic
trajectories and fields, Monte Carlo methods produce stochastic ensembles,
Numerical Optimization supports calibration and simulation-based
optimization, and Scientific Computing provides the software,
parallelization, verification, and reproducibility infrastructure needed
for large computational studies. Simulation also forms a central bridge to
the applied modeling topics of Chapter~13, where biological, physical,
financial, economic, network, and control systems are formulated in
application-specific detail. The next section, High-Performance
Mathematical Computing, develops the algorithms and computational
architectures required to scale mathematical calculations to large
problems, distributed systems, accelerators, and modern high-performance
hardware.
}

%%%%%%%%%%%%%%%%%%%%%%%%%%%%%%%%%%%%%%%%%%%%%%%%%%%%%%%%%%%%
% SECTION SOURCE TRACKING
%%%%%%%%%%%%%%%%%%%%%%%%%%%%%%%%%%%%%%%%%%%%%%%%%%%%%%%%%%%%

%------------------------------------------------------------
% 11.10 Mathematical Simulation
%------------------------------------------------------------
%
% Source basis:
%   - Standard mathematical modeling and simulation methodology.
%   - Standard deterministic and stochastic simulation concepts.
%   - Standard verification, validation, calibration, sensitivity, and
%     uncertainty-quantification principles.
%
% Used for:
%   - mathematical simulation definition;
%   - model versus simulation distinction;
%   - states, parameters, inputs, and outputs;
%   - deterministic and stochastic simulation;
%   - static and dynamic simulation;
%   - continuous- and discrete-time simulation;
%   - time-driven simulation;
%   - event-driven simulation;
%   - event calendars;
%   - queue simulation;
%   - agent-based simulation;
%   - synchronous and asynchronous updates;
%   - compartment models;
%   - conservation diagnostics;
%   - network simulation;
%   - cellular automata;
%   - Markov state transitions;
%   - simulation horizons;
%   - transient and steady-state behavior;
%   - warm-up periods;
%   - finite-horizon simulation;
%   - time-step selection;
%   - physical and numerical time scales;
%   - physical versus numerical instability;
%   - initial-condition sensitivity;
%   - chaotic behavior;
%   - conservation, positivity, and boundedness diagnostics;
%   - stochastic replications;
%   - ensemble means;
%   - simulation confidence intervals;
%   - aleatoric and parametric uncertainty;
%   - uncertainty propagation;
%   - output distributions;
%   - scenario analysis;
%   - parameter sweeps;
%   - local and global sensitivity;
%   - normalized sensitivity;
%   - variance-based sensitivity preview;
%   - calibration;
%   - weighted least squares;
%   - simulation-based parameter estimation;
%   - inverse problems;
%   - structural and practical identifiability;
%   - calibration versus validation;
%   - code and solution verification;
%   - benchmark problems;
%   - convergence studies;
%   - synthetic data;
%   - residual analysis;
%   - simulation error decomposition;
%   - model-form error;
%   - uncertainty quantification;
%   - confidence and prediction intervals;
%   - rare-event simulation;
%   - simulation-based optimization;
%   - noisy objectives;
%   - common random numbers;
%   - surrogate models;
%   - response surfaces;
%   - reduced-order simulation;
%   - digital experiments;
%   - factorial simulation designs;
%   - replication and randomization;
%   - simulation metadata and data management;
%   - checkpointing and restart;
%   - parallel simulation;
%   - parallel random streams;
%   - simulation visualization;
%   - phase portraits;
%   - ensemble visualization;
%   - projectile, predator--prey, epidemic, inventory, reliability, and
%     financial simulation examples;
%   - applications and exercises.
%
% Notes:
%   - High-Performance Mathematical Computing is developed next in
%     Section 11.11.
%   - Mathematical Modeling is developed in Chapter 13.
%   - Numerical ODEs appear in Section 11.4.
%   - Numerical PDEs appear in Section 11.5.
%   - Monte Carlo Methods appear in Section 11.6.
%   - Scientific Computing appears in Section 11.7.
%
%%%%%%%%%%%%%%%%%%%%%%%%%%%%%%%%%%%%%%%%%%%%%%%%%%%%%%%%%%%%